\section{Part F. Reconstruction of schemes and the inertia comparison}

\subsection{Two different conditions that the prompt calls pinned}

The manuscript locus is cote 149, manuscript p.~25v, scan~52. The
transcription says ``1-fid\`ele'' and then asks about full faithfulness
after universal homeomorphisms have been inverted. The actual scan is
not in the supplied image selection; we make no palaeographic assertion
about it. We distinguish a \emph{relative} geometric morphism, which
includes a specified commutation isomorphism over $B G_k$, from a
\emph{closed-point-preserving} relative morphism. Carlson--Haine--Wolf
call the latter \emph{pinned}, in Definition~2.19 of the supplied
preprint \citep{F_CHW2024}. This is an additional condition. LM5-F.5
states it correctly; the definition in R27 of the present prompt omits
it. We use $\operatorname{Hom}^{\mathrm{cp}}_{B G_k}$ for that
restricted Hom-groupoid below, so that no change of definition is hidden.

\begin{claim}{F.1}{Refuted}
With the literal definition of ``pinned'' in R27, F1 is false. More
strongly, there is no characteristic-zero field over which all relative
geometric morphisms between the small \emph{\'etale} topoi of finite-type
schemes are induced by scheme morphisms, even after absolute weak
normalization. This failure already occurs for source $\Spec k$ and
target $\mathbb P^1_k$.
\end{claim}
\begin{proof}
Choose separable closures compatibly with $k(t)\subset k((t))$.
In equal characteristic zero the decomposition group $D_0$ of the
$t$-adic valuation in $G_{k(t)}$ has an exact sequence
\[
1\longrightarrow\widehat\Z(1)\longrightarrow D_0
 \longrightarrow G_k\longrightarrow1.
\]
It has a continuous splitting: in a maximal unramified extension act on
the coefficients, and extend that action by fixing a compatible system
of roots $t^{1/n}$. Every finite extension over the maximal unramified
field is tame, and its tame part is obtained by these roots; hence this
description defines the action on a separable closure, not only on a
finite list of extensions. Restriction to a separable closure of
$k(t)$ gives a continuous homomorphism $s:G_k\to G_{k(t)}$ whose
composite with $G_{k(t)}\to G_k$ is the identity. Equivalently one can
make this construction in the henselization of $k(t)$ at $0$.

A continuous homomorphism of profinite groups induces the geometric
morphism of their continuous-action topoi with inverse image the
restriction functor. The composite
\[
B G_k\xrightarrow{B s} B G_{k(t)}
 =\Spec(k(t))_{\mathrm{et}}
 \longrightarrow (\mathbb P^1_k)_{\mathrm{et}}
\]
is therefore a relative geometric morphism, with its specified
commutation isomorphism over $B G_k$. Its map on the spaces of points
sends the sole point to the generic point of $\mathbb P^1_k$.
Every $k$-morphism $\Spec k\to\mathbb P^1_k$ has closed image.
Isomorphic geometric morphisms induce the same map on subterminal
objects and hence on the underlying sober spaces, so the displayed
morphism is not induced by a scheme morphism. A field is already
absolutely weakly normal in characteristic zero; universal
homeomorphisms preserve the topological obstruction. This is precisely
the example of \citet[Example~2.30]{F_CHW2024}, with the splitting
construction made explicit. The example is valid for $k=\Q$, for
$\Q^{\mathrm{ab}}$, and for $\Q^{\mathrm{solv}}$.
\end{proof}

\begin{claim}{F.2}{Proved}
There is, over every characteristic-zero field $k$, a relative
endomorphism of $(\mathbb G_{m,k})_{\mathrm{et}}$ which induces the
identity on its arithmetic fundamental group and has underlying image
the generic point. It is not induced by an algebraic morphism and does
not factor through $B G_k$. It is, however, constant on underlying
topological points. Consequently it does not settle the
closed-point-preserving, nonconstant test intended in F2.
\end{claim}
\begin{proof}
Put $X=\mathbb G_{m,k}$ and $E=\pi_1(X)$, retaining the full arithmetic
quotient. The map $D_0\to E$ induced by the generic point is an
isomorphism: it is the identity on $G_k$ and the local Kummer inertia
maps isomorphically onto $\pi_1(X_{\bar k})=\widehat\Z(1)$.
Finite \emph{\'etale} covers of $X$ give the canonical geometric
morphism $X_{\mathrm{et}}\to B E$; its inverse image sends a continuous
$E$-set to the corresponding locally constant sheaf. This functor
preserves finite limits and all colimits, and has the usual right
adjoint. Compose it with
\[
B E\simeq B D_0\longrightarrow B G_{k(t)}
 \longrightarrow X_{\mathrm{et}}.
\]
On $\pi_1$ the composite is $E\simeq D_0\to G_{k(t)}\to E$, the
identity. The middle field topos has one underlying point, and the last
map takes it to the generic point. An algebraic endomorphism of $X$ is
$t\mapsto c t^m$, $c\in k^*$ and $m\in\Z$; none has this point
map. Factoring through $B G_k$ would kill the nontrivial geometric
kernel on $\pi_1$, which the displayed morphism does not do. The last
sentence of the claim records the distinction between this fact and
nonconstancy of a point function.
\end{proof}

\subsection{The repaired cyclotomic statement}

We now audit the intended mathematical content of LM5-F.47 and F.50
with the missing closed-point condition restored. The following proof
does not use either AI return as a theorem. It identifies the arithmetic
inputs and supplies the steps from them to the reconstruction criterion.
For an abelian group $A$ write $\widehat A=\varprojlim_n A/nA$ and
$T A=\varprojlim_n A[n]$, using multiplication transition maps in
the second limit. For a field $L$ put
$D(L^*)=\bigcap_{n\ge1}(L^*)^n$.

\begin{claim}{F.3}{Known, with derivation}
Let $k$ be a subfield of a finitely generated extension $M$ of
$\Q^{\mathrm{ab}}$. For every finite extension $L/k$ in a fixed
algebraic closure, there is an overfield $E\supset L$ finitely
generated over $\Q^{\mathrm{ab}}$ such that
\[
 E^*/\mu_\infty\text{ is free abelian},\qquad
 D(E^*)=\mu_\infty,
 \qquad \bigcap_{n\ge1}n A(L)=0
\]
for every abelian variety $A/L$. In particular
$D(L^*)\subset\mu_\infty$.
\end{claim}
\begin{proof}
We quote exactly the forms of the three arithmetic inputs used here.
Amoroso--Zannier, Theorem~1.1, states that for a number field $F$, an
abelian extension $B/F$, and a nonzero algebraic number $a$ which is not
a root of unity, its logarithmic height has a positive lower bound
depending only on $F$ and $[B(a):B]$ \citep{F_AZ2000}. In particular
it is uniformly positive when the latter degree is bounded. Tsujimura,
Proposition~4.9 of RIMS-1957, states that if $F$ is a number field and
$f\in F$, then the field $B$ obtained by adjoining all roots of $f$
has $B^*/B^{*\infty}$ free abelian \citep{F_Tsu2022}. With $f=1$
this gives the needed statement over $F\Q^{\mathrm{ab}}$.
Tsujimura, Theorem~3.1, states that for a number field $F$ and an
abelian variety $A/F$ the natural Kummer map
\[
A(F\Q^{\mathrm{ab}})\longrightarrow
H^1\bigl(G_{F\Q^{\mathrm{ab}}},T A\bigr)
\]
is injective \citep{F_Tsu2020}. Its proof factors it as the natural
map to $\varprojlim_n A(F\Q^{\mathrm{ab}})/n$ followed by an
injection. Thus its kernel is precisely the intersection of the
subgroups $nA(F\Q^{\mathrm{ab}})$.

Take $E=ML$, finite over $M$, and let $E_0$ be the algebraic closure
of $\Q^{\mathrm{ab}}$ in $E$. Since $E$ is finitely generated,
$E_0/\Q^{\mathrm{ab}}$ is finite, hence
$E_0=F_0\Q^{\mathrm{ab}}$ for a number field $F_0$. A normal proper
geometrically integral $E_0$-model $Z$ of $E$ has
\[
 E^*/E_0^*\lhook\joinrel\longrightarrow
 \operatorname{Div}(Z).
\]
Indeed, a rational function with zero divisor on a proper normal
variety is a global unit, and the constant field is $E_0$.
The divisor group is free abelian, as is each of its subgroups.
The height theorem shows $D(E_0^*)=\mu_\infty$: arbitrary roots of
a nontorsion element would have positive heights tending to zero.
The quoted freeness theorem now makes $E_0^*/\mu_\infty$ free.
The exact sequence with this group as kernel and $E^*/E_0^*$ as
quotient splits, because the quotient is free. This proves freeness
of $E^*/\mu_\infty$. Alternatively, applying the divisor map to an
infinitely divisible element and to each of its roots shows they all
lie in $E_0$, proving $D(E^*)=\mu_\infty$ directly.

For the abelian-variety assertion first work over $E$. Spread $A/E$
and a point $P\in\bigcap_n nA(E)$ to an abelian scheme
$\mathcal A\to V$ and a section $P_V$, with $V$ a normal integral
variety over $\Q^{\mathrm{ab}}$. A chosen division $Q_n\in A(E)$
is a generic section of the finite \emph{\'etale} cover
$[n]^{-1}(P_V)\to V$. Normality extends this generic section over
the same $V$: its selected component is finite birational and hence
isomorphic to $V$. Therefore every division extends without further
shrinking. At each closed point $v$, the residue field is finite over
$\Q^{\mathrm{ab}}$, and the abelian variety and point descend to a
number field enlarged by finitely many coefficients. The quoted
Kummer theorem gives $P_V(v)=0$. Closed points are dense and
$\mathcal A$ is separated, so $P_V=0$. Finally embed $A(L)$ into
$A(E)$ to prove the assertion over $L$. This is also the content of
Tsujimura, Proposition~1.2(ii), whose specialization proof we have
just written with the common-domain step explicit.
\end{proof}

\begin{claim}{F.4}{Proved}
Let $L$ be a finite extension of a field $k$ as in F.3, let
$b_1,\ldots,b_r\in L$ be distinct, and put
$V=\mathbb A^1_L\setminus\{b_1,\ldots,b_r\}$. Suppose
$\xi\in\varprojlim_n H^1(V,\mu_n)$ has an actual Kummer value in
$L^*$ at every $L$-point of $V$. Then $\xi$ is the Kummer class of
a unit in $\mathcal O(V)^*$. No discrete valuation of $L$ is assumed.
\end{claim}
\begin{proof}
Since $\operatorname{Pic}(V)=0$ and
$\mathcal O(V)^*=L^*\times\prod_j(t-b_j)^\Z$, write
\[
 \xi=\gamma+\sum_j\lambda_j[t-b_j],\qquad
 \gamma\in\widehat{L^*},\quad\lambda_j\in\widehat\Z.
\]
We construct $x,y\in V(L)$ such that
$r_j=(x-b_j)/(y-b_j)$ are multiplicatively independent modulo roots
of unity. Put $F=\Q(b_1,\ldots,b_r)$. Choose a smooth integral
finite-type model $A$ of $F$ over $\Z[1/N]$, containing the $b_j$
and all inverse differences $(b_i-b_j)^{-1}$. This follows by taking
a smooth affine model over $\Q$, adjoining these functions and
shrinking, then spreading its finite presentation over
$\Z[1/N]$ and shrinking the arithmetic base. Its image contains
all but finitely many rational primes. Choose distinct such primes
$p_i>r$ and a generic point of a component of the fibre at each
$p_i$. Regularity and reducedness of that fibre make its local ring
a discrete valuation ring with uniformizer $p_i$; denote the
valuation by $v_i$. The residues of the $b_j$ are pairwise distinct.
Choose $c_i\in\F_{p_i}$ avoiding them. The Chinese remainder
theorem supplies $x,y\in A$ with
\[
x\equiv b_i+p_i\pmod{p_i^2 A},\qquad
y\equiv c_i\pmod{p_i A}\quad(1\le i\le r).
\]
Thus $v_i(r_j)=\delta_{ij}$. In particular $x,y\in V(L)$ and the
$r_j$ are independent modulo roots of unity. An equality proving
dependence in an overfield would be an equality between these same
elements of $F$, and so the independence persists.

Embed $L$ in $E$ as in F.3. The classes of the $r_j$ in the free
abelian group $E^*/\mu_\infty$ have finite support in a basis and
are independent. There are therefore integer-valued characters
$h_i:E^*/\mu_\infty\to\Z$ for which the square matrix
$N_{ij}=h_i(r_j)$ has nonzero determinant. Let $u_x,u_y\in L^*$
realize the two given evaluations. Applying completed characters
to their ratio gives
\[
 N(\lambda_j)_j=(h_i(u_x/u_y))_i\in\Z^r.
\]
The adjugate matrix shows $(\det N)\lambda_j\in\Z$.
The diagonal copy of $\Q$ meets $\widehat\Z$ in $\Z$, as is
checked one prime at a time; hence $\lambda_j\in\Z$. After
subtracting these integral unit classes, an evaluation at any
$L$-point identifies $\gamma$ with an actual constant Kummer class.
This gives the required unit. For $r=0$ the last evaluation alone
is the proof.
\end{proof}

\begin{claim}{F.5}{Proved}
Let $C/L$ be a smooth geometrically connected affine curve, with $L$
as in F.4. Suppose $\xi\in\varprojlim_n H^1(C,\mu_n)$ has an
actual multiplicative Kummer value at every point over every finite
extension of $L$. After a finite extension of $L$, $\xi$ is the
class of a regular unit of $C$.
\end{claim}
\begin{proof}
Let $\overline C$ be the smooth projective completion and
$S=\overline C\setminus C$. A finite extension makes $S$ split and
supplies a point $x_0\in C(L)$. Residues give a vector
$(\lambda_s)_{s\in S}\in\widehat\Z^S$. For a fixed $s$, choose
by Riemann--Roch a finite map
$\pi:\overline C\to\mathbb P^1_L$ with a simple zero at $s$ and
nonzero values at all other points of $S$. Remove its branch values,
$0,\infty$ and the images of $S$, and, after another finite constant
extension, split this finite set. This gives a finite \emph{\'etale}
map $W\to V$, with $W\subset C$ and $V$ a punctured affine line.
The norm of $\xi|_W$ is a class on $V$ to which F.4 applies: above
an $L'$-point its Kummer value is the product of the norms of the
actual values at the finite residual fields in that fibre.

The valuation of a field norm is the sum of residue degrees times
the local valuations. Consequently the residue of this norm class
at $0$ is $\lambda_s$: the only boundary point above $0$ is the
rational point $s$, and all points inside $C$ have zero residue.
This proves $\lambda_s\in\Z$; no ramification multiplier occurs
in the divisor-of-norm formula. Set $D=\sum_s\lambda_s[s]$.
The localization sequence and the Kummer injection
$\operatorname{Pic}(\overline C)/n\hookrightarrow
H^2(\overline C,\mu_n)$ imply
$[D]\in\bigcap_n n\operatorname{Pic}(\overline C)$. Its degree
is zero, and every one of its division classes has degree zero.
The rational point $x_0$ identifies actual degree-zero line bundles
with $J(L)$, so F.3 gives $[D]=0$. Let $g_0$ have divisor $D$.

The class $\xi-[g_0]$ is residue free. At every finite level its
extension from $C$ to $\overline C$ is unique, since
$H^1(\overline C,\mu_n)\to H^1(C,\mu_n)$ is injective; uniqueness
makes those extensions compatible. Taking limits of the Kummer
sequences on $\overline C$ gives
\[
0\longrightarrow\widehat{L^*}\longrightarrow
\varprojlim_n H^1(\overline C,\mu_n)
\longrightarrow T\operatorname{Pic}(\overline C)
\longrightarrow0.
\]
The left system has surjective transition maps, so there is no
derived-limit obstruction. Each coordinate of the last Tate module
is an infinitely divisible degree-zero Picard class, and F.3 kills
it. Thus the remaining class is constant. Evaluation at $x_0$ makes
it the class of an actual element of $L^*$, completing the proof.
\end{proof}

\begin{claim}{F.6}{Proved}
Let $K$ be an algebraic closure of $k$ as in F.3 and let
$B=\Q^{\mathrm{ab}}\subset K$. Let $C/K$ be a smooth projective
connected curve. Suppose a function $q$ on a cofinite subset of
$C(K)$ satisfies, for every $b\in K$,
\[
 q(x)-b=e_b(x)r_b(x),\qquad e_b(x)\in B^*,
 \quad r_b\in K(C)^*,
\]
outside a finite set depending on $b$. Then $q$ agrees with one
rational function on a cofinite subset. If $r_0$ is constant, that
rational function is constant.
\end{claim}
\begin{proof}
Every finitely generated field $N/B$ inside $K$ has algebraic
constant field $N_0=N\cap\overline\Q$ finite over $B$. The
extension $N/N_0$ is regular, so for algebraic numbers $a$,
$[N(a):N]=[N_0(a):N_0]$. These degrees are unbounded: otherwise
$\overline\Q/N_0$ would have bounded finite subextension degrees,
and hence finite degree, contrary to the height lower bound of F.3
applied to $2^{1/n}$. A fixed finite map of a curve over $N$ to
$\mathbb P^1_N$ therefore supplies closed points of arbitrarily
large degree over $N$. They can avoid any fixed finite set of
$K$-points; the algebraic points in that finite set have bounded
degrees, and the other points cannot equal a closed $N$-point.

Assume $r_0$ is nonconstant. Write the representations of $q$ and
$1-q$ as $A(x)r(x)$ and $B(x)s(x)$, with coefficients in $B^*$.
Descend $C,r,s$ to a finitely generated field $N/B$. A nonzero
rational function $a r+b s-1$, $a,b\in B^*$, has zero divisor of
degree bounded independently of $a,b$, by the fixed pole divisor.
At a point of degree larger than this bound the pointwise equality
forces the rational identity $a r+b s=1$. Absorb those two constants
into $r,s$, obtaining $r+s=1$. It follows that
\[
 (A(x)-B(x))r(x)=1-B(x),
\]
so $q(x)=r(x)$ whenever $r(x)\notin B$, away from finitely many
exceptions.

Choose $a\in K\setminus B$ and apply the same argument to $q+a$.
Its rational factor cannot be constant: otherwise the cofinite
image of the nonconstant $r_0$ would lie in a vector space of
dimension at most two over $B$, whereas $[K:B]=\infty$. Obtain
$q(x)=r_a(x)-a$ where $r_a(x)\notin B$. These two outside-$B$ sets
cover the common cofinite domain, since membership of both rational
values in $B$ would put both $q(x)$ and $q(x)+a$ in $B$.
Their intersection is infinite: after adjoining all coefficients
to $N$, a point of sufficiently large degree cannot belong to a
fibre over a value in $B$ of either of the fixed finite maps $r,r_a$.
On the intersection $r=r_a-a$, hence this is a rational identity.
It proves the assertion on the union.

If $r_0$ is a nonzero constant, rescale $q$ so that its cofinite
values lie in $B$. For $b\notin B$, the factor $r_b$ is constant,
because otherwise its cofinite image would lie in $B+Bb$.
Two different values of $q$ would then force first $r_b\in B$
and then $b\in B$ by subtracting their two equations. Thus $q$
is constant. This argument verifies the finite-exception and
transcendental-coefficient versions explicitly; it does not cite the
algebraic-extension version of Tsujimura's quasi-rationality
criterion outside its hypotheses.
\end{proof}

\begin{claim}{F.7}{Proved}
With the closed-point condition restored, LM5-F.50 is valid: if $k$
is any subfield of a finitely generated extension of
$\Q^{\mathrm{ab}}$, then for all finite-type $k$-schemes $X,Y$
the natural map
\[
\operatorname{Hom}_k(X^{\mathrm{awn}},Y)
\longrightarrow
\pi_0\operatorname{Hom}^{\mathrm{cp}}_{B G_k}
 (X_{\mathrm{et}},Y_{\mathrm{et}})
\]
is bijective. In characteristic zero $X^{\mathrm{awn}}$ is the
seminormalization of $X$. The statement includes the corrected
LM5-F.47 by taking transcendence degree zero.
\end{claim}
\begin{proof}
We first prove the needed point-function assertion for a smooth
geometrically connected affine curve $C/L$. A closed-point-preserving
relative morphism
$\phi:C_{\mathrm{et}}\to(\mathbb G_m)_{\mathrm{et}}$ gives a
$G_L$-equivariant function $f:C(\overline L)\to\overline L^*$.
An $L'$-point maps to an $L'$-point for every finite $L'/L$:
composition gives a relative morphism of closed-point Galois topoi,
whose map to $B G_{L'}$ forces the residue field of the selected
image to embed in $L'$. The specified relative isomorphism also
identifies the coefficient sheaves $\mu_n$, so pulling back Kummer
classes and evaluating at points commute.

F.5 applied to the coordinate class gives, after a finite extension,
a unit $g_0$ such that $f(x)/g_0(x)\in\mu_\infty$ for every
geometric point. If the generic image of $\phi$ is closed,
continuity forces its entire point function into that closed point;
after a finite extension the function is constant. Otherwise, for
each $b\in\overline L$, the inverse image of
$\mathbb G_m\setminus\{b\}$ is a dense open $C_b$, with finite
complement. Applying F.5 to the class of $t-b$ on this open gives
a rational function $g_b$, a unit on $C_b$, with
\[
f(x)-b\in\mu_\infty\, g_b(x) \quad(x\in C_b(\overline L)).
\]
F.6 gives a rational $R$ agreeing with $f$ outside a finite set.
For each $b$, $(R-b)/g_b$ has cofinite image in $\mu_\infty$ if
it is nonconstant, an impossibility; thus it is a constant root of
unity. Fix $x$ and choose $b\ne f(x)$. The identity $R-b=c_b g_b$
first shows that $R$ has no pole at $x$. For every
$b\notin\{f(x),R(x)\}$ it then gives
\[
 \frac{f(x)-b}{R(x)-b}\in\mu_\infty.
\]
If $f(x)\ne R(x)$ the left side, as $b$ varies, fills
$\overline L\setminus\{0,1\}$, contradicting the proper subset
$\mu_\infty$. Hence $f=R$ everywhere. It follows that $R$ is a
regular unit. Galois equivariance descends it to $L$.

If $C(L)\ne\varnothing$, its compatible Kummer class equals the
one pulled back by $\phi$ already over $L$. Over the finite
extension used in F.5, $R/g_0$ takes only root-of-unity values,
so it is constant; evaluation at an $L$-point makes its class zero.
The residual finite-level class over $L$ is geometrically zero,
and the low-degree Hochschild--Serre sequence makes it a constant
class. Evaluation at the $L$-point, which splits the constant-class
inclusion, makes that residual class zero. This proves equality at
all levels without an unwarranted descent of a chosen Kummer root.

We next pass to arbitrary dimension. Let $X/k$ be connected regular
affine of finite type, let $x_0\in X(k)$, and let
$\phi:X_{\mathrm{et}}\to(\mathbb G_m)_{\mathrm{et}}$ be
closed-point-preserving and relative. Regularity in characteristic
zero gives smoothness. The algebraic constant field in $k(X)$
embeds in $k$ by evaluation at $x_0$, so $X$ is geometrically
integral. Choose a smooth projective compactification $Y$ with
normal-crossings boundary $D_1,\ldots,D_s$. The strong resolution
and principalization result of \citet[pp.~1--4, ``The result''
and ``Idea of proof'']{F_EH2002} gives this by operations supported
off the given smooth open. Its hypotheses are characteristic zero
and a finite-type embedded variety, both satisfied here.

Write $h=(h_n)_n$ for the pulled-back coordinate Kummer class.
For every boundary prime $D_i$ it has a residue
$r_i\in\widehat\Z$. After a finite extension choose a rational
smooth point $q$ of a geometric component of $D_i$, away from
the other components, and a smooth geometrically integral projective
curve through $q$ and $x_0$, transverse to $D_i$ at $q$. Here is the
precise Bertini input. Stacks, Lemma~33.47.1 says sufficiently high
powers of an ample bundle, with prescribed vanishing on a finite
subscheme, embed its complement; Lemma~33.47.3 makes their general
divisor smooth on that complement; Lemma~37.32.3 makes it
geometrically irreducible when the base locus has codimension at
least two and the linear system has two independent moving
directions \citep{F_StacksBertini}. Surjectivity of the first-jet
evaluation for sufficiently high powers allows independent
differentials at the prescribed points, including transversality.
Intersect these nonempty opens successively, reducing dimension
to one. The fields are infinite, so each such open has a rational
point. This constructs the specified curve and checks the extra
hypotheses, rather than presuming a curve through fixed points.

The curve calculation makes its restricted Kummer class integral;
its residue at $q$ is $r_i$. To check the comparison, in the regular
local ring at $q$ the sole boundary has equation $z=0$.
At each finite level, after an \emph{\'etale} localization, the
Kummer class is represented by $u z^j$: this local ring is regular
and factorial, and the unit factors away from $z$ have zero
boundary residue. The transverse curve pulls $z$
back to a uniformizer, and $u$ to a unit, so the residues agree
modulo every $n$. Finite separable constant extension has
multiplicity one along the divisor. Thus all $r_i\in\Z$.

The integral divisor $D=\sum_i r_iD_i$ has zero Kummer Chern
class modulo every $n$, hence $[D]\in\bigcap_n n\operatorname{Pic}(Y)$.
This intersection vanishes. Indeed the quotient of
$\operatorname{Pic}(Y)$ by its algebraically trivial subgroup is
a subgroup of the finitely generated geometric N\'eron--Severi
group, by \citet[Theorem~7.3]{F_CJLO2021}. That theorem assumes a
smooth projective scheme over an algebraically closed field, and
is applied to $Y_{\bar k}$. Kill the finite torsion of the quotient
by an integer $e$. If $x$ is divisible by all integers in
$\operatorname{Pic}(Y)$, its quotient is zero, and divisions by
$em$ yield divisions of $x$ by $m$ in the algebraically trivial
subgroup. This subgroup injects into the $k$-points of the Picard
abelian variety; F.3 therefore gives $x=0$. The same argument
kills $T\operatorname{Pic}(Y)$, since each coordinate of a compatible
torsion system is infinitely divisible. The possible Brauer
obstruction only restricts the subgroup of Picard-scheme points
represented by actual line bundles; it does not spoil this
injection, which follows from Hilbert 90.

Choose $g\in k(Y)^*$ with divisor $D$. The class $h-[g]$ is
unramified along every prime of $Y$. A finite-level Kummer torsor
unramified in codimension one on a regular scheme extends \emph{\'etale}
over it by purity. The exact form used is Stacks,
Lemma~58.31.2: a finite \emph{\'etale} cover of a dense open,
unramified along the missing prime divisors, extends when the
missing local rings are regular \citep{F_StacksPurity}.
Its extension is unique by normality and
generic injectivity. The limit Kummer sequence and
$T\operatorname{Pic}(Y)=0$ make the remaining class constant.
Evaluation at $x_0$ supplies an actual $c\in k^*$, so
$h=[c g]$ and $cg(x_0)=\phi(x_0)$.

For any $a\in X(\bar k)$ choose a finite extension making it
rational and a curve through $a$ and $x_0$ by the same Bertini
argument. The curve's actual unit $f_C$ and $cg|_C$ have equal
Kummer classes. Their ratio is divisible by every integer in the
unit group. Its divisor on the projective completion is therefore
zero, so it is constant. Its value at $x_0$ is one. It follows that
$cg(a)=\phi(a)$ for every $a$. The case $\dim X=1$ is the curve
argument, and $\dim X=0$ is immediate.

Finally Carlson--Haine--Wolf, Theorem~4.18, for an arbitrary field,
states the equivalence between full reconstruction and its
condition (4.18.5): for every affine regular connected finite-type
source with a rational point, every pinned map to
$(\mathbb G_m)_{\mathrm{et}}$ has the geometric-point function
of a morphism to $\mathbb G_m$ \citep{F_CHW2024}. We have proved
exactly that condition, and for all finite extensions used in its
descent. Applying the theorem gives the bijection stated above.
This final use is essential: equality of point functions for one
map alone would not identify the entire inverse-image functor.
The theorem's proof uses the same assertion on sliced \emph{\'etale}
objects, reconstructs affine coordinate maps, and applies $h$-descent
in the source. Thus neither a fundamental-group isomorphism nor a
point-function calculation is substituted for the full topos claim.
\end{proof}

\begin{claim}{F.8}{Known}
For a characteristic-zero field $k$, the following is an exact
necessary and sufficient criterion for \emph{closed-point-preserving}
\emph{\'etale} reconstruction: for every connected smooth affine
$X/k$ with $X(k)\ne\varnothing$, every closed-point-preserving
relative morphism $X_{\mathrm{et}}\to(\mathbb G_m)_{\mathrm{et}}$
has on $\bar k$-points the values of one regular unit of $X$.
\end{claim}
\begin{proof}
This is the equivalence (4.18.1)$\Leftrightarrow$(4.18.5) of the
opened Carlson--Haine--Wolf theorem, with ``regular'' replaced by
``smooth'' in characteristic zero and
$X^{\mathrm{awn}}=X$ for regular $X$. The forward implication is
restriction to the stated source and target; the reverse implication
is the reconstructed affine-coordinate and $h$-descent argument
specified at the end of F.7. This is a topos-theoretic criterion,
not an intrinsic classification in terms of the field's arithmetic.
\end{proof}

\paragraph{Literature comparison.}
The September~2026 version of Carlson--Stix has been opened as
well. Its Theorem~4.1 proves reconstruction over sub-$p$-adic
fields; Theorem~6.12 adds a relative result for componentwise
universally open schemes over geometrically unibranch noetherian
Nagata bases with sub-$p$-adic function fields
\citep{F_CS2026}. Proposition~3.17 assumes Kummer-faithfulness
for its point-theoretic open-image dichotomy. These hypotheses
do not cover $\Q^{\mathrm{ab}}$ or $\Q^{\mathrm{solv}}$.
Indeed, the algebraic constant field in a finitely generated
extension of $\Q_p$ is finite over $\Q_p$, whereas either of
these fields contains roots of unity of unbounded $p$-power order.
This also corrects R27's attribution: LM5-F.38 actually cites
Carlson--Stix~2026, not Carlson--Haine--Wolf~2024.
No opened source or targeted search located the full statement
of LM5-F.47 or F.50. This records the search outcome and makes
no assertion of priority or of absence from the literature.

\begin{claim}{F.9}{Unresolved}
An intrinsic necessary-and-sufficient arithmetic classification of
all characteristic-zero fields satisfying closed-point-preserving
reconstruction for every finite-type scheme remains unresolved here.
The corrected $\mathbb G_{m,\Q^{\mathrm{solv}}}$ endomorphism
test is proved in F.23 below. F.31 proves all-scheme
reconstruction over finite extensions of $\Q^{\mathrm{solv}}$,
and F.37 enlarges this to every subfield of a finitely generated
extension of $\Q^{\mathrm{solv}}$. These results give a
sufficient field class. F.38--F.41 below add intrinsic necessary
conditions and explicit excluded field families; their converse
is not proved, so the classification remains unresolved.
\end{claim}
The earlier F.3--F.7 proof does not extend directly to
$\Q^{\mathrm{solv}}$. Its difficulty is not lack of a discrete
valuation: F.4 already removes that assumption. That proof uses the
freeness of an ambient multiplicative group modulo $\mu_\infty$
and the absence of divisible abelian-variety points. Over
$\Q^{\mathrm{solv}}$, the multiplicative group is divisible.
Indeed, for $a\in(\Q^{\mathrm{solv}})^*$ and $n>0$, choose
a finite solvable Galois extension $L/\Q$ containing $a$ and
$\mu_n$. Adjoining $n$-th roots of all the conjugates of $a$
gives an abelian extension of $L$ and a Galois solvable extension
of $\Q$. It therefore lies in $\Q^{\mathrm{solv}}$ and supplies
the desired root. In particular the multiplicative Kummer class
of $2$ is zero although $2$ is not a root of unity.
Thus the coefficient-character argument and
the residual-root-of-unity conclusion cannot be applied. An elliptic
curve over $\Q$ also has its entire $2$-primary torsion rational
over $\Q^{\mathrm{solv}}$: the groups
$\operatorname{GL}_2(\Z/2^n)$ are solvable, since their reductions
to $\operatorname{GL}_2(\F_2)\simeq S_3$ have $2$-group kernels.
These points are divisible by every integer in that field and
invalidate F.3's abelian-variety input. Neither failure constructs
a nonalgebraic pinned morphism of the full \emph{\'etale} topos.
A system of finite-cover maps must still account for all open
subtopoi and their partly ramified charts. F.20--F.23 supply that
comparison for the specific rational endomorphism test.
F.24--F.31 then supply a different all-scheme proof by finite
quotient detection, arithmetic equivariance over a number field,
and the semiabelian Tate-Hom theorem; they do not use the failed
abelian-variety input. The generic-valued map
F.2 has constant underlying image and is not a counterexample
to that corrected test. Under R27's literal definition the universal-field
criterion has the empty solution set by F.1, and therefore cannot
simultaneously contain the incorrectly phrased F.50.

F.16 and F.18--F.19 below give further finite-extension constraints:
the induced profinite exponent is an integer, and the complete point
function is an algebraic monomial. F.20--F.23 complete the comparison
on full topoi by finite-quotient detection and normalization of
charts. The higher-genus and higher-dimensional extensions
are proved separately in F.24--F.31. F.17 refutes the literal
relative-only nonconstant assertion.

\subsection{What cyclotomic normalization does and does not prove}

\begin{claim}{F.10}{Known, with explicit derivation}
Let $K$ have characteristic zero, let $X/K$ be smooth of finite type,
let $D$ be a normal-crossings divisor, and put $U=X\setminus D$.
At a geometric point $\bar x$ with branch set $J$, the strict local
covering-germ group is $T^J$, $T=\widehat\Z(1)$. A chosen path
gives a continuous homomorphism
$j_x:T^J\to\pi_1(U_{\bar K})$, well defined up to geometric
conjugacy. The local decomposition group acts on $T^J$ by the
cyclotomic character times the permutation of branches. A compatible
monomial map
\[
 t'_j=u_j\prod_i t_i^{m_{ji}}
\]
induces the integer contact matrix $(m_{ji})$ on the two geometric
local inertia groups. Units change tangential splittings by Kummer
cocycles but do not change this geometric inertia map.

For a family of smooth hyperbolic curves over a smooth finite-type
base, a horizontal puncture group injects into the geometric fibre
and total groups. If $U$ is a hyperbolic polycurve and $X$ is any
smooth proper normal-crossings compactification, the diagonal
$T\to T^J\xrightarrow{j_x}\pi_1(U_{\bar K})$ has nonzero
image. At a rational boundary point every section factoring through
the punctured henselian neighbourhood therefore has nontrivial
cyclotomic centralizer. Global injectivity of the \emph{full} group
$T^J$ is an additional assertion; it does not follow from these facts.
\end{claim}
\begin{proof}
SGA~1, XIII, Proposition~5.2 assumes a regular local scheme and a
normal-crossings divisor given by part of a regular parameter system;
it says a finite cover tamely ramified on the divisor becomes
\emph{\'etale} after adjoining roots of those parameters, with orders
prime to the residue characteristic. Corollary~5.3 then identifies
the tame fundamental group of the punctured strict localization with
$\prod_{\ell\ne p}\Z_\ell(1)^J$ \citep{F_SGA1}. Here $p=0$,
so every prime occurs. On the cover $z_i^n=t_i$, an inertia tuple
acts by $z_i\mapsto\zeta_i z_i$. Conjugating by an element of the
local arithmetic group sends $\zeta_i$ to its cyclotomic power and
permutes the parameter branches. Units have roots in the strict
henselian ring, which proves both the contact-matrix formula and
its independence of those roots on geometric inertia. Paths account
exactly for the geometric conjugacy ambiguity.

For a punctured surface of type $(g,r)$ with
$2g-2+r>0$, $r>0$, its fundamental group is free on
$a_1,b_1,\ldots,a_g,b_g,c_1,\ldots,c_{r-1}$; the last puncture is
the inverse of $\prod_i[a_i,b_i]\prod_j c_j$. Every puncture word
is nontrivial. A cyclic subgroup generated by a nontrivial word in
a free group has its full profinite topology: represent a cyclically
reduced word by an immersed circle in a rose, complete that finite
labelled graph to a finite covering without identifying its vertices,
and use its primitive homology class to map the resulting finite-index
subgroup to $\Z/n$. Taking the core gives finite quotients in which
the word has order divisible by every prescribed $n$. Thus its
completed cyclic subgroup is $\widehat\Z$ and injects.
SGA~1, XII, Theorem~5.1 identifies finite algebraic and analytic
\emph{\'etale} covers of a locally finite-type complex scheme;
Corollary~5.2 identifies its algebraic fundamental group with the
finite-index completion of its topological group. Descent of the
finite-type pair to a finitely generated field and characteristic-zero
geometric base change transport this argument to $\bar K$.
This proves the horizontal statement, including a single puncture,
whose inertia can vanish in abelianization while remaining injected
in the full group. Injection from the geometric fibre to the total
group is Hoshi, Proposition~2.4(i), valid for a hyperbolic curve
over a connected noetherian separated normal characteristic-zero
base \citep{F_Hoshi2014}. The present smooth base satisfies these
hypotheses.

For diagonal nonvanishing, the further input is Hoshi,
Lemma~2.8: a point of a relative hyperbolic curve over the fraction
field of a strictly henselian trait extends over the trait exactly
when its map to the total fundamental group is trivial. Its base
has the same noetherian, separated, normal, characteristic-zero
hypotheses. Induction gives this criterion for a polycurve: first
extend to the preceding stage, then apply the relative-curve result.
In completed branch coordinates at $\bar x$, send every boundary
parameter to $t$ and all other parameters to zero. This defines an
arc $\Spec\bar K[[t]]\to X_{\bar K}$ with generic point in $U$
and closed point $\bar x$. The Kummer covers identify its map on
fundamental groups with $j_x$ composed with the diagonal.
If that map were zero, the criterion would extend the generic
point to $U$. Separatedness of $X$ identifies this extension with
the original arc, contradicting its boundary closed point.
Finally branch permutations fix the diagonal and the cyclotomic
kernel acts trivially on it. Equivariance gives the centralizer
assertion for every locally factoring section, including nonsplit
branches. This is precisely the extra nonvanishing statement of
CK-C8; it is weaker than full local injection.
\end{proof}

\begin{claim}{F.11}{Proved}
The following precise group-theoretic formulation supplies the
injection step of LM5-F.39. Suppose
$1\to F_r\to G\to Q\to1$, $r\ge2$, is an extension of discrete
groups. Then $\widehat{F_r}\to\widehat G$ is injective and
\[
1\longrightarrow\widehat{F_r}\longrightarrow\widehat G
 \longrightarrow\widehat Q\longrightarrow1
\]
is exact. Suppose in addition free abelian local lattices $A,A_B$
map compatibly to $G,Q$, with contact map $M:A\to A_B$, the map
$\widehat A_B\to\widehat Q$ injective, and $\ker M$ either zero
or infinite cyclic whose generator maps to a nontrivial element of
$F_r$. Then $\widehat A\to\widehat G$ is injective. This applies
at each stage of a recursively compatible stable pointed curve
tower with affine hyperbolic fibres and injected base meridian
lattices: the node kernel is generated by the essential vanishing
cycle, and the marked kernel by an essential puncture cycle.
\end{claim}
\begin{proof}
First $Z(\widehat{F_r})=1$. For a finite quotient $F_r\to A_0$
with kernel $L$, choose a prime $\ell\nmid|A_0|$. The cellular
chains of the finite regular rose covering give
\[
0\to H_1(L,\F_\ell)\to\F_\ell[A_0]^r
\to\F_\ell[A_0]\to\F_\ell\to0.
\]
Semisimplicity gives
$H_1(L,\F_\ell)\simeq\F_\ell\oplus
\F_\ell[A_0]^{r-1}$, a faithful $A_0$-module. In the finite
quotient $F_r/[L,L]L^\ell$, the image in $A_0$ of a central
element of $\widehat{F_r}$ must act trivially on this module, so
must be one. This for every $A_0$ proves centrelessness.

The automorphism group of a finitely generated profinite group is
profinite: intersections of its finitely many open subgroups of
bounded index give a characteristic open basis. Therefore conjugation
$G\to\operatorname{Aut}(\widehat{F_r})$ extends to
$\widehat G$. Its composite with
$\widehat{F_r}\to\widehat G$ is the inner-automorphism map,
which is injective. The quotient of $\widehat G$ by the closure
of $F_r$ classifies exactly the finite quotients of $Q$, proving
the asserted exactness.

Smith normal form shows
$\ker\widehat M=\widehat{\ker M}$: the image of $M$ is free,
the sequence onto that image splits, and its inclusion in $A_B$
stays injective after completion. A kernel element of
$\widehat A\to\widehat G$ first maps to zero in
$\widehat A_B$, then lies in $\widehat{\ker M}$. In the cyclic
case F.10's full-topology argument and the first paragraph of this
claim make the latter map injective, so the element is zero.

In a stable curve chart $xy=t_1^{a_1}\cdots t_s^{a_s}$ the normal
lattice has, relative to the base, one kernel direction represented
by the two opposite branch meridians. On a smooth fibre it is the
vanishing cycle. A vanishing cycle bounding an unmarked disc would
cut off an unmarked genus-zero tree attached at one node; summing
special points over that tree contradicts stability. Thus it is
essential. At a marking the corresponding kernel is a puncture
cycle. Compatible root changes restrict these lattices to finite
index; smooth toroidal subdivisions insert essential annular
directions and preserve the same kernel calculation. The hypotheses
that these are the local lattice maps and that the base lattices
already inject are exactly what is required for the induction;
an arbitrary blowup of a base boundary is not included.

For the geometric application, the required full exact sequence
at a smooth tower stage, including its left injection, is also
given directly by Hoshi, Proposition~2.4(i). The base is smooth,
connected, separated, noetherian and normal. Apply the identical
kernel chase to this profinite sequence. Thus the application does
not presume that a merely right-exact homotopy sequence is left
exact, or that arbitrary profinite completion preserves kernels.

We record CK-B5's logically separate group assertion as well.
Every open subgroup $V$ of $\widehat{F_r}$ is the completion of
the corresponding finite-index discrete free subgroup: cores of
finite-index subgroups show that the induced topology is full.
The graph-covering formula gives its rank
$1+[\widehat{F_r}:V](r-1)\ge2$, so the centre calculation above
applies to $V$. The group $\widehat{F_r}$ is torsion free:
if a nonidentity finite-order $z$ survives modulo an open normal
$N$, put $V=N\langle z\rangle$. Its nonzero cyclic quotient
$V/N$ factors through the torsion-free abelianization of $V$,
where the finite-order element $z$ must vanish, a contradiction.
For an open normal $V$, its centralizer injects into the finite
group $\widehat{F_r}/V$ and is therefore trivial. Passing to a
normal subgroup inside an arbitrary open $H$ proves
$C_{\widehat{F_r}}(H)=1$.
This property passes to extensions: project a centralizer element
to the slim quotient, then centralize the open intersection with
the slim kernel. It passes to surjective inverse limits because
the image of an open subgroup under each quotient projection is
open, and all its coordinate centralizers vanish. Rank one is
excluded, since $\widehat\Z$ centralizes itself. This proves the
CK-B5 comparison without converting slimness into an inertia
recognition theorem.
\end{proof}

\begin{claim}{F.12}{Refuted}
Cyclotomic normalization alone does not imply global injection of a
full normal-crossings inertia group, and finite Kummer root changes
cannot repair a nonzero kernel. A finite root stack also changes
the group and can enlarge the local inertia kernel.
\end{claim}
\begin{proof}
Let $U=(\mathbb P^1_{\Q}\setminus\{0,1,\infty\})^2$ and
blow up the point $(0,2)$ of its product compactification. Write
local coordinates $x=uv$, $y-2=v$ on one blowup chart. The two
boundary components at $u=v=0$ are the strict transform and the
exceptional divisor. Their meridians both map to the old $x=0$
meridian; the loop at $y=2$ is null-homotopic in the second factor.
Thus, identifying the injected old cusp group with $T$, the local
map is
\[
T^2\longrightarrow T,\qquad(a,b)\longmapsto a+b,
\]
with kernel $\{(a,-a):a\in T\}$. Both factors have the stated
cyclotomic action. The diagonal survives by multiplication by two.
The discrete matrix and its reductions are checked in
\texttt{code/F\_checks.py}; the displayed calculation proves the
whole profinite statement.

If $A\ne0$ is a closed subgroup of $\widehat\Z^r$ and $L$ is a
finite-index lattice, choose $n$ with $n\widehat\Z^r\subset L$.
Multiplication by $n$ is injective, so
$0\ne nA\subset A\cap L$. A finite root \emph{cover} only passes
to such a lattice and an open global group, so the kernel persists.
For a finite root \emph{stack} with orders $n_i$, the global group
is the quotient by the closed normal group generated by the
$n_i$-th powers of boundary inertia. This follows directly from
the local charts $[\Spec A[z_i]/(z_i^{n_i}-t_i)/\prod_i\mu_{n_i}]$:
its \emph{\'etale} covers are exactly the tame covers whose inertia
orders divide $n_i$. For $(\mathbb A^1,0)$ the original inertia
map is the identity of $T$, whereas its map to the $n$-th root
stack has kernel $nT$. Taking all divisible orders recovers the
original kernel because every finite quotient has finite inertia
orders. These are the distinct cover and stack statements in
LM5-F.22, and neither contradicts F.11.
\end{proof}

\begin{claim}{F.13}{Proved}
Trait-inertia conditions satisfied by a homomorphism do not, by
themselves, make it the homomorphism of a rational point. This holds
both for trivial inertia and for the positive-multiplicity, injective,
finite-index cusp conditions of F149-C3. Likewise the injection of
actual tangent vectors into the completed Kummer torsor in F149-C8
is not surjective, already over $\Q$.
\end{claim}
\begin{proof}
Set $R=\Q[[t]]$, $K=\Q((t))$, and let
$U/R=\mathbb P^1_R\setminus\{0,1,\infty\}$ be constant. Write
$E=\pi_1(U_{\Q})$ for the group of the total good model. The
coefficient splitting identifies
$G_K=T\rtimes G_{\Q}$, and the tangent vector $1$ at the rational
cusp $0$ identifies its normalizer $D_0\subset E$ with
$T\rtimes G_{\Q}$, with cyclotomic action in both products.

For the zero-inertia example compose $G_K\to G_{\Q}$ with this
cuspidal splitting. It kills inertia. If it came from a $K$-point,
the local Kummer formula would force its proper extension to avoid
the boundary; its reduction would be a point of $U(\Q)$. The
resulting section would then be both a point section and a cuspidal
section over $\Q$, contrary to Stix, Theorem~14. That theorem
assumes injectivity for proper hyperbolic curves over every finite
extension; here it follows from the abelian-variety Kummer
injectivity established in F.3. Thus all its hypotheses are met.

For the positive-inertia example choose
$\lambda\in\widehat\Z$ with $\lambda_2=1$ and
$\lambda_\ell=0$ for all odd primes. It is not an ordinary integer:
an integer divisible by every odd prime is zero, incompatible with
its $2$-adic coordinate. Let
$a=\lambda\kappa_{\Q}(2)\in H^1(G_{\Q},T)$ and choose a
continuous cocycle representative. For an integer $\mu>0$ define
\[
h:G_K\longrightarrow D_0\subset E,
\qquad (z,g)\longmapsto(\mu z+a(g),g).
\]
The cocycle identity proves that $h$ is a continuous homomorphism.
It is injective, its inertia is $\mu T$, and its image has index
$\mu$ in $D_0$. Via the good-reduction pullback
$\pi_1(U_K)=E\times_{G_{\Q}}G_K$, it also determines a section
of the generic arithmetic extension.

Suppose the map to $E$ were, up to geometric conjugacy, induced
by a point $x\in U(K)$. Its nontrivial inertia forces its proper
extension to meet a cusp, and the intersection theorem for cusp
groups (Stix, Proposition~11 and Corollary~12) makes that cusp
$0$. Write $x=t^\mu u$, $u\in R^*$. Pulling back all covers
$z^{1/n}$ shows that its residual cocycle is
$\kappa_{\Q}(\bar u)$: the group $1+tR$ is uniquely divisible
in characteristic zero. A conjugation preserving the nonzero open
subgroup $\mu T$ of the cusp inertia belongs to $D_0$, and its
geometric part belongs to $T$; it changes only the cocycle
representative. Applying the completed $2$-adic-place valuation
$\widehat{\Q^*}\to\widehat\Z$ to the residual class gives the
ordinary integer $v_2(\bar u)$ on the point side and $\lambda$
on the constructed side. They cannot agree. This is a
whole-profinite-group obstruction, not a finite-level experimental
claim.

The same class $a$ gives a splitting of $D_0\to G_{\Q}$ not
induced by any actual nonzero tangent vector. Actual vectors form
a $\Q^*$-torsor; after choosing $1$, their images are precisely
$\kappa_{\Q}(\Q^*)$. Their $2$-valuations are integers, excluding
$a$. The finite CRT residues and the semidirect-product law used
above are checked in \texttt{code/F\_checks.py}.
\end{proof}

\begin{claim}{F.14}{Known, with derivation}
For a hyperbolic curve over a characteristic-zero field $k$ which
distinguishes tori from abelian varieties, a nontrivial closed
subgroup $A\subset\pi_1(U_{\bar k})$ is contained in a cusp
inertia group if and only if it is procyclic and essentially
cyclotomically normalized. A continuous section is cuspidal if
and only if its image cyclotomically normalizes a nontrivial
procyclic geometric subgroup. This supplies a recognition theorem
which the local normalization/injection theorem F.10--F.11 does
not supply. The hypothesis holds for every $k$ as in F.3, in
particular every field finitely generated over $\Q$.
\end{claim}
\begin{proof}
Stix, Theorem~20, states the equivalence of containment in inertia,
procyclic cyclotomic normalization, and procyclic essential
cyclotomic normalization, under exactly the displayed field and
hyperbolicity hypotheses. Corollary~21 identifies the maximal
such subgroups, and Corollary~22 gives the section statement
\citep{F_Stix2008}. To spell out the last implication, Theorem~20
places $A$ in $I_c$. If an element of the section image normalizes
$A$, the two conjugate cusp groups meet nontrivially. Proposition~11
makes them equal, and Corollary~12 identifies their arithmetic
normalizer with the cusp decomposition group. The section image
therefore lies in that group. The reverse direction is the local
Kummer action. The section surjects onto $G_k$, so the cusp is
$k$-rational.

To verify the field hypothesis, split a given torus after a finite
extension $k'/k$. A nonzero equivariant map
$T S\to T A$ would give, at some prime $\ell$, a nonzero map
$\Z_\ell(1)\to T_\ell A$. Restriction to
$k'(\mu_\infty)$ gives nonzero invariant compatible
$\ell$-power torsion. Every coordinate is divisible by all
positive integers in $A(k'(\mu_\infty))$, since multiplication by
integers prime to $\ell$ is invertible on the compatible
$\ell$-primary tower. This field is again a subfield of a
finitely generated extension of $\Q^{\mathrm{ab}}$, so F.3
excludes it. Hence $\operatorname{Hom}_{G_k}(T S,T A)=0$.

No part of this argument identifies every noncuspidal section with
a rational point. That would require the section-surjectivity
assertion separately retained in Part D.
\end{proof}

\begin{claim}{F.15}{Proved: exact scope of the eight-row comparison}
None of F149-C1--C6, C8, C9 asserts that an arbitrary section
comes from a rational point solely because of its inertia. C2,
C3, and C5 do contain converses within the class of homomorphisms
already known to come from rational geometric sections. Those
conditional converses are valid and are stronger than the bare
injection statement of LM5-F.39. They do not turn into an
arbitrary-section recognition theorem. The same distinction applies
to C9's tangential specialization data.
\end{claim}
\begin{proof}
Here is the comparison, keeping the quantifiers of the original rows.
The manuscript references are those of the opened row statements.
\begin{longtable}{@{}>{\raggedright\arraybackslash}p{.12\textwidth}>{\raggedright\arraybackslash}p{.39\textwidth}>{\raggedright\arraybackslash}p{.41\textwidth}@{}}
\toprule
Row & Actual assertion and locus & Converse and its disposition\\
\midrule\endhead
C1 & Hartogs extension of a rational section of a smooth pair or
tower; 149, pp.~32, 34--35, 41; scans~65, 70--72, 83.
& No inertia-recognition assertion. Its input is an actual section
on a dense open.\\
C2 & An actual trait point extends iff its map to the total good
model kills inertia; pp.~36--38, scans~74--78.
& A genuine conditional converse. Proper extension and the nonzero
Kummer multiplicity prove it. The arbitrary-section extension is
false by F.13's zero-inertia example.\\
C3 & For an actual point meeting a cusp, inertia is multiplication
by its positive contact order, with the stated index; pp.~37--39,
44, scans~75--79, 90.
& Equivalent conditions for that already geometric point. The
positive-inertia converse for an arbitrary homomorphism is false
by F.13's $\lambda\kappa(2)$ example.\\
C4 & Exact homotopy and reconstruction of a centre-free extension
from its outer action; p.~39, scan~79, also~64, 73.
& No claim about which sections are geometric. Neither exactness
nor centrelessness implies section surjectivity.\\
C5 & An actual rational section extends over a smooth base iff it
kills divisorial inertia; pp.~32, 34--36, 41--42;
scans~65, 70--74, 83--84.
& A genuine conditional converse using C2 and C1. An arbitrary
unramified section can be cuspidal and need not be an interior
point section; F.13 refutes that extension.\\
C6 & Normalizer of a specified cusp group and its Chern-class
extension; pp.~39--41, scans~79--82.
& The subgroup is already specified as cusp inertia. A section
lying in its normalizer is cuspidal by definition. Recognition
from cyclotomic normalization is the extra theorem F.14.\\
C8 & Actual normal vectors inject into the completed Kummer
torsor of splittings; p.~42, scan~85.
& No surjectivity is asserted. Such surjectivity is false over
$\Q$, by F.13. Replacing normal vectors by conormal frames
changes the Kummer sign.\\
C9 & An actual rational section fails to extend on a divisorial
locus and has its leading-term tangential specialization;
pp.~46--47, scans~93--95.
& The rational section is part of the input. An arbitrary
contact-order datum need not be geometric, again by F.13.\\
\bottomrule
\end{longtable}

For completeness the geometric implications underlying this table
can be checked directly. Over a trait a point extends uniquely to
the proper curve. If it meets the boundary with equation
$z=t^\mu u$, the pullback of $z^{1/n}$ gives inertia
$\tau\mapsto\mu\tau$ in the injected cusp group. If it avoids
the boundary the map factors through the trait. These are exactly
the two directions of C2 and the positive-multiplicity equivalences
in C3. For a higher-dimensional regular base, killing divisorial
inertia extends the point at every height-one local ring. Its domain
then has complement of codimension at least two. The extension
step C1 is geometric: after resolving the rational map to a proper
positive-genus curve, exceptional projective-space fibres map
constantly, since no nonconstant map from $\mathbb P^1$ to that
curve exists. Descend each blowup using a finite birational graph
over the normal base. A boundary intersection is an effective
Cartier divisor, so cannot be supported in codimension two. For
genus zero with at least two removed points, trivialize one boundary
section \emph{\'etale}-locally and apply Hartogs to the resulting
regular function; the other boundary intersections are again
Cartier divisors. The extensions glue by separatedness. Induction
on a tower uses its actual Cartesian pullbacks. This proves the
conditional converse in C5 and explains the divisorial statement
in C9. Its leading coefficient transforms under $t\mapsto at$
as $u\mapsto a^{-\mu}u$, giving the stated Kummer change
$-\mu\kappa(a)$.

The sign in C6 is also consistent with F.10. If $N$ is the normal
line, choose frames $e_j=g_{ij}e_i$ and scalar coordinates
$w_i=g_{ij}w_j$. Roots $r_i^n=w_i$ transform by
$r_i=h_{ij}r_j$, $h_{ij}^n=g_{ij}$. Their triple defect
$h_{ij}h_{jk}h_{ki}\in\mu_n$ is $+c_1(N)\bmod n$ with the
positive Kummer convention. The conormal line is $N^\vee$ and has
the opposite sign. Thus this is extra extension-class information,
not a recognition of geometric point sections.

LM5-F.21 concerns the local tube and its possible Chern
transgression; F.22 concerns root stacks and covers; F.23 changes a
cuspidal model to a stable model; F.35 is the abstract closed-subgroup
criterion for extension of a covering germ as a sub-germ; F.39 is
the injection and normalization induction of F.10--F.11. These are
compatible with the eight rows just listed. CK-B5 is a slimness
statement, not an inertia-recognition statement: the centralizer
of an open subgroup vanishes in nonabelian free profinite groups,
and the property passes through extensions and surjective inverse
limits. Rank one is excluded. CK-C8 asserts nonzero diagonal
boundary inertia and a boundary-local centralizer implication;
it explicitly allows the blowup kernel of F.12. Neither it nor
LM5-F.48 asserts its arbitrary-section converse. The distinction
between local injection, cyclotomic recognition in F.14, and point
section surjectivity in Part D is therefore essential in all three
returns.
\end{proof}

\subsection{A finite-extension constraint for the unresolved solvable-field test}

The following additional test uses finite extensions of
$\Q^{\mathrm{solv}}$. It rules out a proposed source of exotic
endomorphisms without asserting reconstruction of the full topos.

\begin{claim}{F.16}{Proved}
Put $k=\Q^{\mathrm{solv}}$. For every finite extension $L/k$,
\[
 D(L^*)=k^*.
\]
Consequently, if $\phi:X_{\mathrm{et}}\to X_{\mathrm{et}}$ is a
closed-point-preserving relative geometric morphism for
$X=\mathbb G_{m,k}$, then its pullback of the canonical compatible
Kummer class is
\[
 \phi^*[t]=m[t]\quad\text{for one }m\in\Z.
\]
For every $\alpha\in\overline\Q^*$ its induced point value satisfies
\[
 \phi(\alpha)=c_\alpha\alpha^m,
 \qquad c_\alpha\in k^*.
\]
In particular a nonintegral profinite exponent cannot give the
endomorphism requested in F2. This claim alone does not prove that $c_\alpha$ is independent of
$\alpha$ or that $m\ne0$ for a nonconstant point map; F.19 below
settles those two point-function issues. This Kummer calculation
alone does not identify the full geometric morphism; F.20--F.23
provide the additional comparison.
\end{claim}
\begin{proof}
We first establish a descent calculation. Enlarge a finite extension
of $k$ to $L=Mk$, where $M/\Q$ is finite Galois, and put
$B=M\cap k$. Both $M/B$ and $k/B$ are Galois and their intersection
is $B$, so restriction identifies
\[
 \operatorname{Gal}(L/B)=
 \operatorname{Gal}(M/B)\times\operatorname{Gal}(k/B).
\]
Write $G=\operatorname{Gal}(L/k)$ for the first factor. Given a
finite list of elements of $L$, choose a number field $S\subset k$
containing $B$ such that the list lies in the number field $F=MS$.
Suppose $z^n=a\in F^*$ with $z\in L^*$. For $\sigma\in G$ and
$\tau\in\operatorname{Gal}(L/F)$ we have
$\tau z/z\in\mu_n\subset k$. The two factors commute, and
$\sigma$ fixes these roots of unity. Hence
\[
 \tau\left(\frac{\sigma z}{z}\right)
 =\frac{\sigma(\tau z/z)}{\tau z/z}
       \frac{\sigma z}{z}
 =\frac{\sigma z}{z}.
\]
It follows that $\sigma z/z\in F$ and
\begin{equation}\label{eq:F16:radical_descent}
 \frac{\sigma a}{a}\in(F^*)^n.
\end{equation}
The important point is that $F$ can be fixed before the root $z$
and the integer $n$ are chosen.

A number field satisfies $D(F^*)=\{1\}$. Indeed the valuations of
an element in this intersection all vanish, and Dirichlet's unit
theorem then reduces the assertion to a finitely generated abelian
group, whose intersection of all power subgroups is trivial. If
$a\in D(L^*)$, choose $F$ containing $a$ and apply
\eqref{eq:F16:radical_descent} for every $n$. Thus
$\sigma a/a=1$ for every $\sigma\in G$, so $a\in k^*$.
Conversely $k^*$ is divisible by the radical argument in F.9.
This proves $D(L^*)=k^*$ for the enlarged Galois extension, and
then for every intermediate finite extension as well. In particular
a proper finite extension of $k$ is not radically closed: any of its
elements outside $k$ fails some power-divisibility test.

The finite Kummer sequence and $\operatorname{Pic}(X)=0$ give
\[
 H^1(X,\mu_n)=\mathcal O(X)^*/\mathcal O(X)^{*n}
 =\Z/n,
\]
because $\mathcal O(X)^*=k^*\times t^\Z$ and $k^*$ is divisible.
The relative rigidification identifies the pullback of $\mu_n$ with
$\mu_n$, compatibly in $n$. We may therefore write
$\phi^*[t]=\lambda[t]$ with $\lambda\in\widehat\Z$.
For a finite $L/k$ and $\alpha\in L^*$, the composite
$\Spec(L)_{\mathrm{et}}\to X_{\mathrm{et}}
 \xrightarrow{\phi}X_{\mathrm{et}}$ is pinned. The elementary
field-source reconstruction statement of
\citet[Section~3, proposition on reconstruction for fields]{F_CHW2024}
identifies it with a unique $\beta\in L^*$: it factors through its
closed image point, whose residue field is finite separable over
$k$, and the relative map of field topoi gives its embedding in $L$.
Functoriality of Kummer theory gives
\begin{equation}\label{eq:F16:point_kummer}
 \kappa_L(\beta)=\lambda\kappa_L(\alpha).
\end{equation}

There exist a finite Galois $M/\Q$, an element
$\sigma\in\operatorname{Gal}(Mk/k)$, and $\alpha\in M^*$ such
that $\sigma\alpha/\alpha$ is not a root of unity. Here is a
concrete existence argument requiring no assumption on finite
extensions of $k$. The polynomial $T^5-6T+3$ is Eisenstein at $3$.
Its derivative has two real zeros $\pm(6/5)^{1/4}$, with positive
and negative polynomial values, respectively. Thus it has exactly
three real roots, and complex conjugation is a transposition in its
transitive degree-five Galois group. A transitive group of prime
degree is primitive; the graph of conjugates of that transposition
is connected, so those transpositions generate $S_5$. Its splitting
field $M$ therefore has group $S_5$, and
$\operatorname{Gal}(Mk/k)=A_5$. Choose $\sigma\in A_5$ moving a
root $\theta$. Since $M$ contains only finitely many roots of unity,
all but finitely many integers $N$ have
$\alpha=\theta+N\ne0$ and $\sigma\alpha/\alpha$ not a root of
unity: for each root of unity $\zeta\ne1$, the equation
$\sigma\theta+N=\zeta(\theta+N)$ excludes at most one $N$,
and $\zeta=1$ excludes none.

For this $L=Mk$ and $\alpha$, let $\beta=\phi(\alpha)$, and
choose the fixed number field $F=MS$ containing $\alpha$ and
$\beta$. If $a_n\in\Z$ represents $\lambda$ modulo $n$,
\eqref{eq:F16:point_kummer} says
$\beta\alpha^{-a_n}=z_n^n$ for some $z_n\in L^*$. Applying the
descent calculation gives
\[
 \frac{\sigma\beta}{\beta}
 \left(\frac{\sigma\alpha}{\alpha}\right)^{-a_n}
       \in(F^*)^n \quad\text{for every }n.
\]
The group $F^*/\mu(F)$ is free abelian: its unit subgroup modulo
torsion is free of finite rank, and its quotient is a subgroup of
the free group of fractional ideals, so that exact sequence splits.
Choose a homomorphism $h:F^*/\mu(F)\to\Z$ which is nonzero on
$\sigma\alpha/\alpha$. The displayed relations imply
\[
 h(\sigma\beta/\beta)
  =\lambda h(\sigma\alpha/\alpha)\quad\text{in }\widehat\Z.
\]
Since both character values are integers and the second is nonzero,
$\lambda\in\Q\cap\widehat\Z=\Z$. Write it as $m$.
For an arbitrary finite $L/k$ and $\alpha\in L^*$,
\eqref{eq:F16:point_kummer} now gives
$\beta\alpha^{-m}\in D(L^*)=k^*$, which proves the stated
point formula. The remaining distinctions in the last sentence of
the claim are essential: a scalar calculation in the finite-\emph{\'etale}
subtopos supplies neither the maps on all open subtopoi nor their
compatibility.
\end{proof}

\begin{claim}{F.17}{Refuted}
With R27's literal relative-only definition of pinned, the
nonconstant endomorphism assertion in F2 is false over every
characteristic-zero field. There is a relative geometric endomorphism
of $(\mathbb G_{m,k})_{\mathrm{et}}$ whose point map sends $1$ to
$1$ and every other point to the generic point. It is nonconstant
on points and is not algebraic. It is not closed-point preserving,
so this does not refute F2 with CHW's actual definition restored.
\end{claim}
\begin{proof}
Let $X=\mathbb G_{m,k}$, $Z=\{1\}$, $U=X\setminus Z$, and
\[
 T=\Spec\mathcal O^h_{X,1},\qquad K_h=\operatorname{Frac}
       (\mathcal O^h_{X,1}),\qquad D=G_{K_h}.
\]
Write $j:U\hookrightarrow X$ and $i:Z\hookrightarrow X$.
If $I\subset D$ is inertia, then $I=\widehat\Z(1)$ and
$D/I=G_k$. The coordinate $u=t-1$ defines
$\Spec K_h\to\mathbb G_{m,u}$ and an isomorphism
\[
 D\xrightarrow{\sim}\pi_1(\mathbb G_{m,u}).
\]
Indeed it induces the identity on $G_k$ and the Kummer
isomorphism on the geometric kernels; all ramification of the
strict henselian characteristic-zero trait is tame. Composing the
map induced by $u:U\to\mathbb G_{m,u}$ with this identification
gives
\[
 q:U_{\mathrm{et}}\longrightarrow BD
\]
over $BG_k$. Its restriction to the local decomposition group at
$1$ is the identity $D\to D$.

For a continuous discrete $D$-set $B$, the nearby-stalk calculation
therefore gives a natural $G_k$-equivariant isomorphism
\begin{equation}\label{eq:F17:nearby}
 i^*j_*q^*B\simeq B^I.
\end{equation}
The geometric stalk on the left is obtained by taking sections
over the punctured strict henselization at $1$, hence by taking
invariants under its local inertia. The preceding decomposition
map identifies both that action and the residual $G_k$-action
with the ones on the right. This calculation applies to every
continuous $D$-set: such a set is a disjoint union of finite
$D$-orbits, and the punctured strict henselization is connected.
Thus $q^*B$ is a union of finite locally constant sheaves; no
common \emph{\'etale} trivialization of an arbitrary infinite
union is assumed.

Artin gluing identifies $X_{\mathrm{et}}$ with the category of
triples
\[
 (A,C,v),\qquad A\in BG_k,\quad C\in U_{\mathrm{et}},
 \quad v:A\longrightarrow i^*j_*C.
\]
For $T_{\mathrm{et}}$ the corresponding triples are
$(A,B,v:A\to B^I)$, with $B\in BD$. These are statements about
all sheaves of sets; the exact gluing theorem and its
henselian-trait form are recalled in
\citet[pp.~1--2]{F_Cornut2024}, citing SGA~4, IV, 9.5.
Define a functor on these full triple categories by
\[
 \Phi^*(A,B,v)=
 (A,q^*B,\ A\xrightarrow{v}B^I
                  \xrightarrow{\sim}i^*j_*q^*B).
\]
Finite limits and all colimits of the triple categories are
computed on their two components, with the induced connecting
map; the gluing functors preserve finite limits. Consequently
$\Phi^*$ preserves finite limits and colimits. For completeness
its right adjoint can be given explicitly. For $(A',C,v')$ put
$B=q_*C$, use the natural map
\[
 B^I\simeq i^*j_*q^*q_*C\longrightarrow i^*j_*C,
\]
and take $A=A'\times_{i^*j_*C}B^I$, with the projection
$A\to B^I$. The adjunction follows directly from
$q^*\dashv q_*$ and the fibre-product condition. Thus $\Phi^*$
is the inverse image of a geometric morphism
$\Phi:X_{\mathrm{et}}\to T_{\mathrm{et}}$. On both components
it commutes with the structural inverse image from $BG_k$;
\eqref{eq:F17:nearby} respects that commutation. Hence $\Phi$
is a relative geometric morphism with the required isomorphism.

Compose $\Phi$ with the morphism $T_{\mathrm{et}}\to
X_{\mathrm{et}}$ induced by the henselization. The generic open
of $T$ pulls back to $U$, and its closed stratum pulls back to
$Z$, with identity on $BG_k$. The point map of the composite is
therefore exactly the one in the claim. Its two values show
nonconstancy. A morphism of finite-type $k$-schemes sends closed
points to closed points, whereas the displayed morphism sends
the rational point $2$ to the generic point. It is consequently
nonalgebraic. Finally $X$ is smooth, so absolute weak
normalization does not alter this example.
\end{proof}


\begin{claim}{F.18}{Proved}
Put $k=\Q^{\mathrm{solv}}$. Finite extensions of $k$ have
unbounded degrees, even when required to contain a prescribed finite
extension $L/k$. Let $U\subset\mathbb G_{m,L}$ be a nonempty
open, and let
$\xi\in\varprojlim_n H^1(U,\mu_n)$. Suppose that every evaluation
of $\xi$ at every point over every finite extension $E/L$ is the
compatible Kummer class of an actual element $b_x\in E^*$.
Then $\xi$ is the compatible Kummer class of an actual unit
$h\in\mathcal O(U)^*$, and
\[
 b_x/h(x)\in k^*
\]
for every such point. The statement applies to arbitrary finite
$L/k$, although these fields are not in general radically closed.
\end{claim}
\begin{proof}
There is a proper finite extension of $k$, by the $S_5$ polynomial
in F.16. If $E/k$ is any proper finite extension, choose
$a\in E\setminus k$. Since $D(E^*)=k^*$, some $n$-th root of
$a$ does not belong to $E$. Adjoining it gives a proper finite
extension of $E$. Iterating gives degrees at least doubling.
The construction can start with any prescribed proper finite $L/k$;
when $L=k$, first use the displayed proper extension. The primitive
element theorem then supplies points of arbitrarily large degree
over $L$. This argument uses F.16's divisibility calculation and
not an unproved assertion about faithful Kummer lines.

Write $P_0=t$ and let $P_1,\ldots,P_r\in L[t]$ be the distinct
monic irreducible polynomials for the other removed closed points.
Since this coordinate ring is a localization of a principal ideal
domain, its Picard group is zero and
\[
 \mathcal O(U)^*=L^*\times\prod_{j=0}^rP_j(t)^\Z.
\]
Choose $x_0\in U(L)$, and an actual $b_0\in L^*$ representing
the evaluation there. The Kummer sequence expresses $\xi$ as
\begin{equation}\label{eq:F18:normalized_class}
 \xi=\kappa_L(b_0)+\sum_{j=0}^r\lambda_j
       \bigl([P_j(t)]-[P_j(x_0)]\bigr),
 \qquad \lambda_j\in\widehat\Z.
\end{equation}
Thus a possibly nontrivial arithmetic constant class over $L$
has been removed by one actual evaluation, not assumed to vanish.

Take a proper finite extension $L'/L$ as above. Choose a finite
Galois number field $M/\Q$ containing generators of $L$ and $L'$
over $k$, all coefficients and roots of the $P_j$, and $x_0,b_0$.
Then $L'\subset Mk$, and we can choose a nonidentity
$\sigma\in\operatorname{Gal}(Mk/L)$. Its restriction to $M$ is
nonidentity; it fixes the coefficients of every $P_j$, but need
not fix their roots.

There are infinitely many rational primes splitting completely
in $M$. One elementary justification suffices here. For a primitive
integral element of the Galois field $M$, infinitely many primes
divide nonzero values of its minimal polynomial. Indeed, if the
possible primes were $p_1,\ldots,p_s$, choose $f(a)\ne0$ and
apply this claim to the nonconstant integer polynomial
\[
 \frac{f(a+f(a)p_1\cdots p_s T)}{f(a)},
\]
which is congruent to $1$ modulo every $p_i$: an integer value of
absolute value greater than $1$ has a new prime divisor. Outside
the finitely many index and discriminant primes, a linear factor
of the minimal polynomial gives a residue-degree-one prime in
$M$, which forces complete splitting because $M/\Q$ is Galois.
Choose distinct completely split primes $p_i$, $0\le i\le r$,
avoiding coefficient denominators, discriminants and pairwise
resultants of the $P_j$, and with $p_i>\sum_j\deg P_j$.

Choose a place $v_i$ of $M$ over $p_i$ and a root $r_i\in M$
of $P_i$. Complete splitting and $\sigma\ne1$ imply
$v_i\ne\sigma^{-1}v_i$. All these places for different $i$
lie over distinct rational primes. In the ring of integers of
$M$ with the finitely many required denominators inverted, the
Chinese remainder theorem provides $\alpha$ satisfying
\[
 \alpha\equiv r_i+p_i\pmod{\mathfrak p_{v_i}^2},
 \qquad
 \alpha\equiv c_i\pmod{\mathfrak p_{\sigma^{-1}v_i}},
\]
where $c_i\in\F_{p_i}$ avoids all roots of all the $P_j$.
The residue-size bound ensures such $c_i$ exist. Taylor expansion
at the simple root and the fixed-coefficient identity give
\[
 v_i(P_j(\alpha))=\delta_{ij},\qquad
 v_i(P_j(\sigma\alpha))
 =v_{\sigma^{-1}v_i}(P_j(\alpha))=0.
\]
In particular $\alpha\in U(Mk)$, and the ratios
\[
 R_j=P_j(\sigma\alpha)/P_j(\alpha)
 \quad\text{satisfy}\quad v_i(R_j)=-\delta_{ij}.
\]
The second congruence is imposed at $\sigma^{-1}v_i$ precisely
because $\sigma$ may permute the roots of $P_j$.

Let $b_\alpha\in Mk$ realize the evaluation of $\xi$. Choose
one fixed number field $F=MS\subset Mk$ containing $b_\alpha$,
where $S\subset k$ is finite and contains $M\cap k$.
For integer representatives $a_{j,n}$ of $\lambda_j$ modulo $n$,
\eqref{eq:F18:normalized_class} gives $z_n\in Mk$ with
\[
 \frac{b_\alpha}{b_0}
 \prod_j\left(\frac{P_j(\alpha)}{P_j(x_0)}\right)^{-a_{j,n}}
 =z_n^n.
\]
Apply F.16's fixed-number-field descent to $\sigma$. Since
$\sigma$ fixes $b_0$ and every $P_j(x_0)$, it yields
\[
 \frac{\sigma b_\alpha}{b_\alpha}
       \prod_jR_j^{-a_{j,n}}\in(F^*)^n.
\]
The integer-valued characters
$h_i=v_i\circ N_{F/M}:F^*\to\Z$ satisfy
$h_i(R_j)=-[F:M]\delta_{ij}$. Consequently
\[
 [F:M]\lambda_i=-h_i(\sigma b_\alpha/b_\alpha)\in\Z.
\]
As $\Q\cap\widehat\Z=\Z$, every $\lambda_i$ is an integer.
The normalized expression for $\xi$ is now the Kummer class of
\[
 h(t)=b_0\prod_j\left(\frac{P_j(t)}{P_j(x_0)}\right)^{\lambda_j}
       \in\mathcal O(U)^*.
\]
At any finite-extension point, the quotient $b_x/h(x)$ has zero
compatible Kummer class. F.16 identifies that kernel with $k^*$,
proving the last assertion.
\end{proof}

\begin{claim}{F.19}{Proved}
Let $k=\Q^{\mathrm{solv}}$. Every closed-point-preserving relative
geometric endomorphism $\phi$ of $(\mathbb G_{m,k})_{\mathrm{et}}$
has, on all geometric closed points, the form
\[
 \phi(\alpha)=c\alpha^m,\qquad
 c\in k^*,\quad m\in\Z.
\]
If $m=0$, the entire geometric morphism is the algebraic constant
map. Thus a nonconstant pinned endomorphism has $m\ne0$ and
agrees with an algebraic monomial on all geometric closed points.
These point-function calculations alone do not identify the
nonconstant \emph{full-topos morphism}; F.20--F.23 below complete
that additional step.
\end{claim}
\begin{proof}
Let $f$ denote the induced $\overline\Q$-point function. F.16
already gives an integer $m$ and
$f(\alpha)=c_\alpha\alpha^m$ with $c_\alpha\in k^*$.
Suppose first that $m=0$. Choose $\beta\notin k$ and put
$L=k(\beta)$. Every geometric closed-point value of $\phi_L$
lies in $k$, so none is $\beta$. The inverse image of the
closed point $\beta$ is closed in $\mathbb G_{m,L}$. Any
nonempty closed subset of this Jacobson scheme contains a closed
point. Thus this inverse image is empty; the generic point cannot
be an isolated exception. The morphism $\phi_L$ therefore factors
through the open complement of $\beta$. Pull back its unit
$t-\beta$ and apply F.18 on all of $\mathbb G_{m,L}$. The
resulting unit has the form $a t^r$, $a\in L^*$, $r\in\Z$,
and
\[
 f(\alpha)-\beta=d_\alpha a\alpha^r,
 \qquad d_\alpha\in k^*.
\]
The left-hand side lies in $L$. If $r\ne0$, a point of degree
$[L(\alpha):L]>|r|$ supplied by F.18 would therefore satisfy
$\alpha^r\in L$, a contradiction. Hence $r=0$. Comparing with
$\alpha=1$ gives
\[
 \frac{f(\alpha)-\beta}{f(1)-\beta}\in k^*.
\]
Since both $f(\alpha)$ and $f(1)$ belong to $k$ while $\beta$
does not, this equation forces $f(\alpha)=f(1)$. The inverse
image of that closed rational point is closed and contains all
closed points of the source, hence is the whole source. By the
closed-subtopos factorization, $\phi$ factors through that
$\Spec(k)_{\mathrm{et}}$. Relativity over $BG_k$ identifies this
factorization with the ordinary algebraic constant morphism.

Now assume $m\ne0$. The point map is nonconstant: otherwise
$\phi$ would factor through a closed point and its geometric
Kummer exponent would be zero. For each $b\in k^*$ the inverse
image of $b$ is a proper closed subset of $\mathbb G_m$, hence
finite. In particular its complement $U_b$ is a nonempty open.
Pull back $t-b$ on the target complement of $b$. F.18 provides
$h_b\in\mathcal O(U_b)^*$ with
\[
 f(\alpha)-b=d_{\alpha,b}h_b(\alpha),
 \qquad d_{\alpha,b}\in k^*,\quad \alpha\in U_b(\overline\Q).
\]
For $b=1$ this gives
\[
 c_\alpha\alpha^m-d_{\alpha,1}h_1(\alpha)-1=0.
\]
After clearing the fixed denominators of $t^m$ and $h_1(t)$,
the numerator degree is bounded independently of the two scalar
values. Choose $\alpha$ outside the finite exceptional fibre
with degree over $k$ greater than this bound. A nonzero polynomial
of that bounded degree cannot vanish at $\alpha$. We therefore
obtain an identity of rational functions, and hence
\[
 h_1(t)=A t^m+B,\qquad A,B\in k^*.
\]
If $\alpha\notin k$, then $\alpha^m\notin k$, since $k$ is
radically closed and contains every root of unity. Such an
$\alpha$ cannot satisfy $f(\alpha)=1$ either. Comparing the
coefficients of $1$ and $\alpha^m$ in the preceding evaluation
equation gives
\[
 c_\alpha=-A/B=:c,\qquad f(\alpha)=c\alpha^m
       \quad(\alpha\notin k).
\]
For any $b\in k^*$, apply the same bounded-degree argument to
$h_b$ at a point of sufficiently large degree, where this
monomial formula is already known. It gives
\[
 h_b(t)=e_b(c t^m-b),\qquad e_b\in k^*.
\]
Every zero of $c t^m-b$ in $\mathbb G_m$ must lie in the removed
fibre $f^{-1}(b)$, because $h_b$ is a unit on $U_b$. For any
$x\in k^*$ choose $b=cx^m$. The zero at $x$ then forces
$f(x)=b$. This proves the monomial formula at every geometric
closed point, including the rational ones.

The last distinction in the claim is substantive. The proposition
on reconstruction from points in Section~4 of
\citet{F_CHW2024} requires algebraic realizations after pullback
to \emph{every} separated quasicompact \emph{\'etale} target
chart. Such charts and their pullbacks include covers of punctured
$\mathbb G_m$ with arbitrary genus. F.18 uses a rational source
with zero Picard group; it supplies no vanishing theorem for the
Picard Tate modules of these other curves. Thus this point-function argument alone does not establish equality
of full-topos morphisms. F.20--F.23 use finite-quotient detection
and normalization of charts to complete that comparison without
asserting vanishing of those Picard Tate modules.
\end{proof}

\subsection{Completion of the corrected solvable-field endomorphism test}

\begin{claim}{F.20}{Proved}
Put $k=\Q^{\mathrm{solv}}$, and let
$U\subset\mathbb A^1_k$ be a nonempty open. For every finite
quotient $\pi_1(U)\twoheadrightarrow G$ and every $n\ge5$, there
is a closed point $x\in U$ of degree $n$ over $k$ whose
arithmetic decomposition group maps onto all of $G$:
\[
 G_{k(x)}\longrightarrow\pi_1(U)\longrightarrow G
 \quad\text{is surjective}.
\]
This detects arithmetic finite quotients, including their constant
quotients. It does not assert that $k$ is Hilbertian or that
$k$-rational points suffice.
\end{claim}
\begin{proof}
Let $Y\to U$ be the connected finite \emph{\'etale} $G$-torsor
corresponding to the quotient. Descend the equations, the entire
$G$-action, and the torsor isomorphism to a finite Galois solvable
number field $F/\Q$ contained in $k$. The descended cover is
connected, since its base change to $k$ is connected. Retain the
notation $Y\to U$ for it, and let $M$ be the algebraic closure of
$F$ in $F(Y)$. There is an exact sequence
\[
 1\longrightarrow N\longrightarrow G\longrightarrow
 P:=\operatorname{Gal}(M/F)\longrightarrow1,
\]
where $N$ is the geometric group. Connectedness after base change
implies $M\cap k=F$; both extensions of $F$ are Galois.

We have $F^{\mathrm{solv}}=k$. Indeed, the normal closure over
$\Q$ of a finite solvable Galois extension of $F$ is solvable:
its group over $F$ embeds into a product of the solvable groups
of its conjugates, and its quotient over $\Q$ is
$\operatorname{Gal}(F/\Q)$. Conversely every element of $k$
lies in a finite solvable Galois extension of $F$. Thus every
solvable quotient of $P$ is trivial, since its fixed subfield
would lie in $M\cap k=F$. In particular $P$ is perfect.

Form the ordered $n$-fold fibre product of $Y$ over its common
constant field $M$, and remove the loci where two of its images
in $U$ coincide. Call the resulting scheme $T$. The factors are
geometrically integral over $M$, so $T$ is geometrically integral
over $M$ and connected over $F$. Over the unordered configuration
space
\[
 B=\operatorname{Sym}^n(U)_{\mathrm{dist}}
\]
it is a connected finite \emph{\'etale} Galois cover with group
\begin{equation}\label{eq:F20:common_constant_wreath}
 H=\bigl(G\times_P\cdots\times_PG\bigr)\rtimes S_n.
\end{equation}
The coordinate automorphisms must induce the same automorphism of
$M$; this is the fibre-product condition. Permutation of the
ordered factors supplies the literal subgroup $S_n$. Together
these actions are free and transitive on the geometric fibres
over an unordered configuration. The source $T$ has \emph{not}
been quotiented by $S_n$: that quotient would have smaller degree
and would not in general be $H$-Galois.

Since $U\subset\mathbb A^1_F$, the coefficients of the monic
polynomial of a configuration identify $B$ with a nonempty open
of $\mathbb A^n_F$. Hilbert irreducibility over the number field
$F$ gives a connected specialization with Galois group $H$, say
$\Omega/F$. Precisely the needed statements are
\citet[Proposition~3.3.1 and Theorem~3.4.1, pp.~23, 25]{F_Serre2007}.
The specialization proposition assumes an irreducible Galois
cover; it does not require its total space to be geometrically
irreducible over $F$. Thus the nontrivial constant field $M$
causes no extra hypothesis here.

Let $H^\infty$ be the stable term of the derived series of $H$.
Since $k=F^{\mathrm{solv}}$,
\[
 \Omega\cap k=\Omega^{H^\infty},\qquad
 \operatorname{Gal}(\Omega k/k)=H^\infty.
\]
We verify that a coordinate stabilizer in $H^\infty$ projects
onto the whole corresponding copy of $G$. First, the literal
$A_n$ is perfect and hence is contained in $H^\infty$. The image
of $H^\infty$ in the permutation group is exactly $A_n$.
Second, $H^\infty\to P$ is onto, because the image of every
derived term is the corresponding derived term of the perfect
group $P$. Multiplication by the inverse literal permutation
therefore shows that
\[
 B_0=H^\infty\cap\ker(H\to S_n)
\]
still surjects onto $P$. Finally $[N^n,A_n]\subset B_0$, and
its projection to any selected coordinate contains all of $N$:
put $a\in N$ in that coordinate and commute it with an even
permutation moving the coordinate elsewhere. The chosen
coordinate of the resulting balanced tuple is $a$ or $a^{-1}$.
It follows that the projection $B_0\to G$ contains $N$ and
surjects onto $G/N=P$, hence is surjective.

The permutation image $A_n$ is transitive. The specialized
configuration consequently gives over $k$ one closed point $x$
of degree $n$. Its decomposition group inside $\Omega k/k$ is
the stabilizer of one coordinate. The preceding calculation
shows that this stabilizer surjects onto the Galois group of the
corresponding $G$-torsor fibre. Equivalently the original
$G_{k(x)}\to G$ is onto, as required.
\end{proof}

\begin{claim}{F.21}{Proved}
Let $k=\Q^{\mathrm{solv}}$, let $X=Y=\mathbb G_{m,k}$, and
let $f(t)=ct^m$ with $c\in k^*$ and $m\ne0$. If a pinned
relative geometric morphism
$\phi:X_{\mathrm{et}}\to Y_{\mathrm{et}}$ agrees with $f$ on
all geometric closed points, then their generic finite
\emph{\'etale} pullback functors are naturally isomorphic over
$BG_k$.
\end{claim}
\begin{proof}
The generic point maps to the generic point: otherwise its image
would be closed, and continuity would make the whole point map
constant. Hence for every nonempty open $V\subset Y$ the
inverse-image open is
\[
 \phi^{-1}(V)=f^{-1}(V)=:U,
\]
again an open rational curve. A finite separable algebra over
$k(Y)$ extends to a finite \emph{\'etale} cover over some $V$.
Pull it back along $\phi|_U$ and take its generic fibre.
Shrinking $V$ does not change that fibre, so this defines the
generic pullback functor. It preserves finite limits, finite
coproducts and covers, and hence defines a continuous generic
Galois homomorphism. This construction does not presuppose an
unproved factorization theorem for arbitrary generic-valued
topos morphisms.

Choose generic fibre functors compatibly with the specified
identification on the base field. The two maps are continuous
homomorphisms over $G_k$,
\[
 h_\phi,h_f:G_{k(X)}\longrightarrow G_{k(Y)}.
\]
Let $q:G_{k(Y)}\twoheadrightarrow Q$ be any finite quotient.
Choose $V$ over which its cover is finite \emph{\'etale}; both
$q h_\phi$ and $q h_f$ factor through $\pi_1(U)$. Form their
\emph{joint} image
\[
 J=\operatorname{im}\bigl((q h_\phi,q h_f):
                   \pi_1(U)\to Q\times Q\bigr).
\]
F.20 supplies a closed point $x\in U$ whose decomposition group
surjects onto $J$.

The two point-source morphisms from
$\Spec k(x)_{\mathrm{et}}$ to $Y_{\mathrm{et}}$ are isomorphic
as relative geometric morphisms. Indeed each factors through its
closed image. The residue field of that image is finite separable
over $k$, and its topos is the corresponding slice of $BG_k$.
Relative morphisms into that slice correspond to actual
$k$-embeddings of the residue field into $k(x)$; equality of
geometric point values identifies these embeddings. This is also
the field-source reconstruction statement in Section~3 of
\citet{F_CHW2024}.

It follows that the restrictions of the two homomorphisms to the
entire decomposition group are conjugate by one element. They
are not merely elementwise conjugate. Choosing the paths to
$x$ over the same geometric point of $\Spec k$ makes the
conjugator geometric. Let $a_Q$ be its image in $Q$.
Surjectivity of the decomposition group onto the joint image $J$
then gives
\[
 q h_\phi=\operatorname{Int}(a_Q)\,q h_f
\]
on the whole generic source group.

Put $\Gamma^{\mathrm{geom}}=G_{\bar k(Y)}$ and define
\[
 C_Q=\{a\in\Gamma^{\mathrm{geom}}:
        q h_\phi=\operatorname{Int}(q(a))\,q h_f\}.
\]
These are nonempty closed subsets of a compact profinite group.
The quotient by the intersection of finitely many kernels refines
all the corresponding quotients; applying the same joint-image
argument to it gives their finite-intersection property.
Consequently there is one $a\in\bigcap_Q C_Q$. It conjugates
the two generic homomorphisms on all finite quotients. Since $a$
is geometric, the induced natural isomorphism of generic
pullback functors respects the identification over $BG_k$.
\end{proof}

\begin{claim}{F.22}{Proved}
Let $X,Y$ be smooth connected separated finite-type curves over a
characteristic-zero field, and let $f:X\to Y$ be finite
\emph{\'etale}. Suppose a geometric morphism
$\psi:X_{\mathrm{et}}\to Y_{\mathrm{et}}$ takes the generic
point to the generic point and is equipped with a natural
isomorphism of its generic finite \emph{\'etale} pullback functor
with that of $f$. Then $\psi\simeq f_{\mathrm{et}}$ on the full
small \emph{\'etale} topoi. If the given generic comparison is
relative to an identified base-field topos, the full comparison
is relative as well.
\end{claim}
\begin{proof}
We give the step from generic covers to all charts explicitly.
By \citet[Theorem~A.8 and Corollary~A.9]{F_CHW2024}, the inverse
image under $\psi$ of a separated quasicompact \emph{\'etale}
chart is represented by a separated quasicompact \emph{\'etale}
scheme. The source is topologically noetherian and the target
is qcqs, exactly as required there. Its proof uses the
complemented diagonal to obtain a separated \emph{\'etale}
algebraic space, which is a scheme by quasi-finiteness; no
additional representability assumption is being made.

A separated \emph{\'etale} chart over a normal connected curve
is canonically an open in the normalization in its finite generic
algebra \citep[Expos\'e~I, Proposition~10.1]{F_SGA1}.
Quasicompactness leaves finitely many generic components, and
normalization is finite because the curves are excellent. Let
$A$ be a finite separable $k(Y)$-algebra, let $S\to Y$ be its
finite normalization, and put $Z=S^{\mathrm{et}}$. The
\emph{\'etale} base change $f^*S$ is normal, so it is the
normalization of $X$ in $k(X)\otimes_{k(Y)}A$. Moreover
\begin{equation}\label{eq:F22:etale_locus}
 (f^*S)^{\mathrm{et}}=f^*Z,
\end{equation}
because being \emph{\'etale} is local for the \emph{\'etale}
topology on the base. The prescribed generic comparison
therefore identifies $\psi^*Z$ with an open in $f^*S$.
Since $\psi^*Z$ is \emph{\'etale} over $X$, it lies in the
locus in \eqref{eq:F22:etale_locus}, giving a canonical inclusion
\begin{equation}\label{eq:F22:open_inclusion}
 \psi^*Z\hookrightarrow f^*Z.
\end{equation}

First this inclusion forces equality of the underlying point
maps. Suppose a closed $x\in X$ has
$\psi(x)=y'\ne y=f(x)$. If $y'$ is closed, choose
$h\in k(Y)^*$ with $v_y(h)=1$ and $v_{y'}(h)=0$; if it is
generic, only the first condition is needed. Approximation for
these discrete valuations gives such $h$. Explicitly, one can
use a common affine neighbourhood and the Chinese remainder
theorem, or prescribe a first jet at $y$ and a nonzero value at
$y'$ by Riemann--Roch on a smooth proper completion. The quadratic
extension $k(Y)(\sqrt h)$ is totally ramified at $y$, and is
unramified at $y'$ when the latter is closed, since $h$ is then
a unit and $2$ is invertible. Its maximal \emph{\'etale} locus
$Z$ has empty geometric fibre at $y$ and nonempty fibre at
$y'$ (including the generic case). Inclusion
\eqref{eq:F22:open_inclusion} makes the fibre of $\psi^*Z$ at
$x$ empty, whereas the inverse-image stalk formula makes it
nonempty. This contradiction proves $\psi(x)=f(x)$.
The generic image is already prescribed, so the point maps agree.
In particular a closed-point-preservation assumption is not
needed for this rigidity lemma.

Now take any finite normalization $S$ and its maximal
\emph{\'etale} locus $Z$. At every geometric closed point
$\bar x$, equality of the underlying image points gives
\[
 |(\psi^*Z)_{\bar x}|=|Z_{\psi(\bar x)}|
     =|Z_{f(\bar x)}|=|(f^*Z)_{\bar x}|.
\]
The fibres are finite, and the inclusion in
\eqref{eq:F22:open_inclusion} is therefore onto on them.
Generic fibres already agree. Thus that inclusion is an
isomorphism, canonically extending the prescribed generic one.
Use it to form the sliced geometric morphism
\[
 \psi_Z:(f^*Z)_{\mathrm{et}}\longrightarrow Z_{\mathrm{et}}.
\]
Its generic functor is identified with that of the finite
\emph{\'etale} algebraic map $f_Z:f^*Z\to Z$.
Apply the ramification argument just proved to these curves.
It shows that the underlying point maps of $\psi_Z$ and $f_Z$
agree. If the curves are disconnected, their finitely many
clopen components are assigned according to their generic
supports, and the argument applies on each relevant connected
component. This uses only the already proved point-map step,
not the desired full comparison recursively.

Every open subobject $V\subset Z$ consequently has the same
inverse-image open under $\psi_Z$ and $f_Z$. Every separated
quasicompact \emph{\'etale} chart $V\to Y$ is such an open in
the maximal \emph{\'etale} locus of its finite normalization.
We have therefore obtained $\psi^*V\simeq f^*V$ on every such
chart, extending the fixed generic comparison. This includes
partly ramified completions: ramified and unramified points may
lie over the same base point, and the maximal \emph{\'etale}
locus followed by its open subobjects selects exactly the latter
points and any desired subset of them. We have not assumed that
finite covers of Zariski opens alone generate every chart.

For a morphism between two charts, the two maps induced by these
identifications agree on generic fibres. Their source is reduced
with dense generic points and their target is separated over
$X$. The equalizer is therefore the whole source, so the
identifications are natural. Separated quasicompact
\emph{\'etale} charts form a basis for the small
\emph{\'etale} topology \citep[Lemma~A.4]{F_CHW2024};
inverse-image functors preserve colimits, so this natural
isomorphism extends to all sheaves. Finally it agrees with the
specified base-field comparison on finite constant objects by
generic uniqueness. Every continuous discrete $G_k$-set is a
coproduct of finite orbits. Thus the full isomorphism is relative
when the generic one is.
\end{proof}

\begin{claim}{F.23}{Proved}
With the actual closed-point-preserving meaning of pinned, the
endomorphism test in F2 has an affirmative answer:
\[
 \operatorname{End}_{k}(\mathbb G_m)
 \xrightarrow{\sim}
 \operatorname{Hom}^{\mathrm{cp}}_{BG_k}
       ((\mathbb G_m)_{\mathrm{et}},(\mathbb G_m)_{\mathrm{et}}),
 \qquad k=\Q^{\mathrm{solv}},
\]
where the right-hand side denotes isomorphism classes with their
specified relative structures. Every such endomorphism is
induced by exactly one $t\mapsto ct^m$, $c\in k^*$,
$m\in\Z$, and it is nonconstant exactly when $m\ne0$.
This endomorphism theorem alone does not imply all-scheme
reconstruction; that property is proved separately in F.31.
\end{claim}
\begin{proof}
F.19 proves the entire constant case and identifies the point
function in the nonconstant case with $f(t)=ct^m$, $m\ne0$.
This algebraic map is finite \emph{\'etale} in characteristic
zero. F.21 gives the relative generic comparison, and F.22
extends it to a relative isomorphism of the full small
\emph{\'etale}-topos morphisms. Uniqueness follows from the
point functions, equivalently from the ordinary description of
the units of $k[t,t^{-1}]$. Both the rational-source hypothesis
in F.18--F.19 and the rational configuration space in F.20 have
been used; a general source has not been substituted without
proof. The literal relative-only statement remains refuted by
F.17. Thus the two definitions have different answers even for
this specific endomorphism test.
\end{proof}

\subsection{Detection beyond rational sources}

\begin{claim}{F.24}{Proved}
Let $k=\Q^{\mathrm{solv}}$, let $L/k$ be finite, and let
$U/L$ be a smooth separated geometrically connected curve. Every finite
quotient $\pi_1(U)\twoheadrightarrow G$ is detected by a closed
point: there exists $x\in U$ for which
\[
 G_{L(x)}\longrightarrow\pi_1(U)\longrightarrow G
\]
is surjective. Such detecting points have unbounded degrees over
$L$. In particular the rational-source restriction in F.20 is
not necessary for finite-quotient detection. No Hilbertianity
assertion about $L$ is made.
\end{claim}
\begin{proof}
Let $Y\to U$ be the connected finite \emph{\'etale}
$G$-torsor, and let $C$ be the smooth proper completion of $U$.
Descend $C,U,Y$, the entire $G$-action, and the torsor
isomorphism to a number field $F\subset L$. Enlarge $F$ inside
$L$ until $L=Fk$. The extension $L/F$ is Galois and
prosolvable, since it is a base change of $k/\Q$.
The model $C/F$ is geometrically connected. The descended $Y$
is connected because $Y_L$ is connected. If $M$ is the
algebraic closure of $F$ in $F(Y)$, there is an exact sequence
\[
 1\longrightarrow N\longrightarrow G\longrightarrow
 P=\operatorname{Gal}(M/F)\longrightarrow1,
\]
with $N$ the geometric group. The curve $Y$ is geometrically
connected over $M$, and connectedness of $Y_L$ gives
$M\cap L=F$. Here $P$ need not be perfect. Put $B=C\setminus U$.
Normalization extends the cover to a finite map
$\bar Y\to C$ of smooth proper curves, ramified only over $B$.

Choose a line bundle $A/F$ on $C$ of degree
$n\ge\max(5,2g(C)+3)$. A suitable pair of $F$-rational
sections defines a degree-$n$ morphism
\[
 r:C\longrightarrow\mathbb P^1_F
\]
with only simple ramification, distinct branch values,
no ramification at $B$, distinct images of the geometric points of $B$,
and $r(B)$ disjoint from all branch values.
For completeness, these generic-pencil assertions follow from
jet counts. For every length-at-most-four subscheme $Z\subset C$
we have $H^1(C,A(-Z))=0$, so restriction of sections to those
jets is onto. A base point imposes two equations with one varying
point parameter. Ramification order at least three likewise
imposes two equations on the local ratio jet. Two distinct
ramification points with the same value impose three equations,
with two point parameters; their diagonal limit is covered by
the preceding higher-ramification condition. Ramification at a
fixed boundary point, or equality of two fixed boundary images,
imposes one equation. A ramification point sharing the image of
a fixed boundary point imposes two equations with one varying
point parameter. The diagonal limit is excluded by
unramifiedness at that boundary point. Thus every bad incidence
has smaller dimension than $H^0(C,A)^2$. The complement of
their closures contains a nonempty open defined over $F$,
which has $F$-rational points because $F$ is infinite.
Base-point-freeness gives $r^*\mathcal O(1)=A$ and hence degree
$n$. Tensor powers of an ample line bundle supply unbounded
such degrees; a curve's index may restrict the available integers.

Let $V\subset\mathbb P^1_F$ be the complement of the branch
values and $r(B)$. Let $H$ be the arithmetic Galois group of the
Galois closure of $F(\bar Y)/F(\mathbb P^1)$. Label the $n$
blocks over a geometric point of $V$, and choose a point in each
$G$-torsor fibre inducing the same embedding of $M$. Monodromy
then embeds
\begin{equation}\label{eq:F24:wreath_embedding}
 H\subseteq (G\times_P\cdots\times_PG)\rtimes S_n.
\end{equation}
Every coordinate induces the same automorphism of $M$.
Restriction $H\to P$ is surjective. The full constant field of
the Galois closure may exceed $M$; no equality is needed.
The geometric monodromy $H^{\mathrm{geom}}$ lies in
$N^n\rtimes S_n$ and acts transitively on the $n|N|$ points
of one geometric $M$-component. That action is faithful:
monodromy commutes with the global deck action of $G$, whose
translates exhaust the other components.

Embed $\bar F$ into $\C$ and apply SGA~1, XII, Theorem~5.1
and Corollary~5.2 \citep{F_SGA1}. Small loops around the
removed values generate the geometric monodromy. Their actions
have exactly two possible forms. At a simple branch value of
$r$, the action is a decorated transposition $r_{ij}(a)$:
it swaps blocks $i,j$ with mutually inverse labels
$a,a^{-1}\in N$ and fixes the others. It has order two,
since $\bar Y\to C$ is unramified over that value.
At the image of a boundary point the action is a
single-coordinate element $e_i(b)$, $b\in N$.
The pencil conditions ensure that all other blocks are
unramified there. Conjugating local loops to the chosen base
point preserves these two forms.

The transposition graph on the $n$ blocks is connected, since
the projection is the transitive monodromy of the geometrically
connected curve $C$. Choose a spanning tree. Conjugating by
one element of $N^n$, chosen successively along that tree,
makes all its edge labels one. Its edge transpositions now
generate the literal $S_n\subset H^{\mathrm{geom}}$.
Let $N_0\subset N$ be the subgroup generated by all remaining
reflection labels and single-coordinate labels. Every generator
has coordinates in $N_0$, so the orbit of $(1,i)$ lies in
$N_0\times\{1,\ldots,n\}$. Transitivity on the chosen
component forces $N_0=N$. Multiplying a remaining reflection
by its literal transposition gives
\[
 d_{ij}(a)=e_i(a)e_j(a^{-1})
          \in H^{\mathrm{geom}}\cap N^n.
\]
Together with the single-coordinate generators, these provide
base elements of support at most two whose coordinate labels
generate $N$. We do not assert that the unrestricted group
$N^n$ lies in geometric monodromy; an Abel--Jacobi constraint
can prevent that stronger assertion.

Let $H^\infty$ be the stable derived term of $H$. The literal
$A_n$ is perfect for $n\ge5$, hence lies in $H^\infty$.
For one of the preceding base elements $u$ and a selected
coordinate $i$ in its support, choose $\tau\in A_n$ with
$\tau^{-1}(i)$ outside that support. Then
\[
 [u,\tau]\in H^\infty\cap N^n,\qquad
 [u,\tau]_i=u_i u_{\tau^{-1}(i)}^{-1}=u_i.
\]
Normality of $H^\infty$ and conjugation by the literal $S_n$
move these labels to every coordinate. Consequently
$H^\infty\cap N^n$ projects onto $N$ at every coordinate.
This is the balanced-label assertion needed below, even in the
presence of the Abel--Jacobi constraint.

Hilbert irreducibility over $F$ gives a specialization
$\Omega/F$ with full group $H$, by
\citet[Proposition~3.3.1 and Theorem~3.4.1]{F_Serre2007}.
As in F.20, irreducibility over $F$ suffices; geometric
irreducibility is not required. Place $\Omega$ in $\bar\Q$
with its constant copy of $M$ equal to the chosen $M$, and set
\[
 J=\operatorname{Gal}(\Omega L/L)
   =\operatorname{Gal}(\Omega/(\Omega\cap L))\subset H.
\]
Since $L/F$ is prosolvable, $H/J$ is solvable, so
$H^\infty\subset J$. Since $M\cap L=F$, restriction
$J\to P$ is onto.

Fix a block $i$ and put $J_i=J\cap\operatorname{Stab}(i)$.
Its coordinate projection to $G$ is a homomorphism and its
image contains $N$ by the commutator calculation. It also
surjects onto $P$: for $h\in J$ with a prescribed $P$-image,
choose a literal $\tau\in A_n$ moving $h(i)$ back to $i$.
Then $\tau h\in J_i$ has the same $P$-image. Thus $J_i\to G$
is surjective. This uses only transitivity of $A_n$; it does
not assume $P$ perfect or require an odd permutation in $J$.

Finally $A_n\subset J$ makes the specialized degree-$n$
fibre of $C\to\mathbb P^1$ over $L$ transitive, hence one
closed point $x$ of degree $n$. It lies in $U$ because the
specialization avoids $r(B)$. The image of $G_{L(x)}$ in $J$
is precisely $J_i$, and its action on $Y_x$ is the coordinate
projection $J_i\to G$. The latter is onto, proving the claim.
\end{proof}

\begin{claim}{F.25}{Proved}
Let $L/\Q^{\mathrm{solv}}$ be finite, and let $X,Y/L$ be
smooth separated geometrically connected curves. Suppose a pinned
relative geometric morphism $\phi:X_{\mathrm{et}}\to
Y_{\mathrm{et}}$ and a nonconstant algebraic map $f:X\to Y$
have the same geometric closed-point values. Their generic
finite \emph{\'etale} pullback functors are naturally isomorphic
over $BG_L$. If $f$ is finite \emph{\'etale}, then
$\phi\simeq f_{\mathrm{et}}$ on the full small topoi.
\end{claim}
\begin{proof}
Continuity forces the generic image to be generic. For every
nonempty target open $V$, the inverse-image open is the same
$U=f^{-1}(V)$ for both morphisms. The generic functors are
therefore defined as in F.21. For a finite generic target
quotient, extend its cover over $V$ and take the joint image
of the two maps from $\pi_1(U)$. F.24 detects that entire
finite arithmetic joint image at one closed point. Field-source
reconstruction supplies a single relative conjugator on its
whole decomposition group, hence on the joint image. The same
compactness argument in the geometric kernel as in F.21 gives
one relative generic comparison. If $f$ is finite
\emph{\'etale}, F.22 extends it over all charts and sheaves.
\end{proof}

The remaining obstruction in this $\Q^{\mathrm{solv}}$ approach
is consequently not a failure of closed-point detection on
higher-genus curves. F.24 detects
all finite quotients, and F.25 compares morphisms once an
algebraic point function is available. The next step, supplied in F.26--F.28 below, is to construct
that point function for an arbitrary pinned morphism
$C_{\mathrm{et}}\to(\mathbb G_m)_{\mathrm{et}}$.
The rational-source Kummer calculation F.18 used
$\operatorname{Pic}(C)=0$; a general curve can have a nonzero
Picard Tate module. Finite-quotient detection does not, by
itself, turn pointwise actual Kummer values into one rational
unit. F.26--F.28 add arithmetic equivariance, the semiabelian
Tate-Hom theorem, and norm-residue integrality. No vanishing of
the whole Picard Tate module or intrinsic field classification
is asserted.

\subsection{Point-theoretic classes on curves: removal of the Picard obstruction}

\begin{claim}{F.26}{Proved}
Let $k=\Q^{\mathrm{solv}}$, let $L/k$ be finite, and let
$U/L$ be a smooth separated geometrically connected curve with
$x_0\in U(L)$. Suppose
$\xi\in\varprojlim_n H^1(U,\mu_n)$ has actual multiplicative
Kummer values at every point over every finite extension of $L$.
After a finite constant extension splitting the boundary, choose
a basis $u_1,\ldots,u_s$ of the regular units modulo constants,
normalized by $u_j(x_0)=1$. Then
\[
 \xi=\kappa(b_0)+\sum_{j=1}^s\lambda_j\kappa(u_j),
 \qquad b_0\in L^*,\quad\lambda_j\in\widehat\Z,
\]
where $b_0$ is the given value at $x_0$. Thus the Picard Tate
component of this particular point-theoretic class vanishes.
For a proper curve the class is constant. No vanishing of the
whole Picard Tate module is asserted.
\end{claim}
\begin{proof}
Write $C$ for the smooth proper completion and $D=C\setminus U$.
After the indicated finite extension, descend $C,D,U,x_0$ to
a number field $F\subset L$ with $L=Fk$, with all boundary
points rational over $F$. Let $A=J_D$ be the reduced-modulus
generalized Jacobian and $j:U\to A$ its Abel--Jacobi map,
normalized at $x_0$. Its universal property identifies
\[
 \operatorname{Hom}(A,\mathbb G_m)
 \simeq\{u\in\mathcal O(U)^*:u(x_0)=1\}.
\]
Its Tate module is the geometric abelian fundamental group.
These facts are the curve case of
\citet[Section~3 and Propositions~4.1--4.3]{F_SpiessSzamuely2003}.
For clarity, the general source has a finite kernel in its
abelian fundamental-group comparison; that kernel is killed by
homomorphisms to $\Z_\ell(1)$. For curves the N\'eron--Severi
quotient in that statement has no torsion: it is $\Z$ in the
proper case and zero after a geometric boundary point is present.
Evaluation along $j$ is the usual point Kummer cocycle, as is
seen directly by pulling back the multiplication-by-$\ell^r$
covers of $A$.

Fix a prime $\ell$, put $H=G_L$, and set
\[
 V=V_\ell A,qquad
 W=V_H=V/\langle(h-1)v:h\in H,\ v\in V\rangle_{\Q_\ell}.
\]
The subspace is closed because $V$ is finite dimensional.
Since $H$ is normal in $G_F$, $W$ is a $G_F$-representation
on which $H$ acts trivially. The image $W_0$ of $T_\ell A$
is a $G_F$-stable full $\Z_\ell$-lattice.
For a point $x\in U(E)$, $E/L$ finite, project the Kummer
cocycle of $j(x)$ to obtain
\[
 c_x:G_E\longrightarrow W_0.
\]
It is a homomorphism independent of division-point choices,
since $G_E\subset H$ acts trivially on $W$. Its
$\Q_\ell$-span is unchanged by a finite extension of $E$:
restricting a homomorphism to a finite-index subgroup changes
its image by finite index.

The images of all $c_x$ span $W$. Otherwise an annihilating
nonzero functional, rescaled to be integral on $W_0$, defines
a nonzero homomorphism $\pi_1(U)\to\Z_\ell$, zero on the
base-point copy of $G_L$ and on every closed-point decomposition
group. Here the rational section splits the arithmetic extension,
and $\mu_\infty\subset L$ makes the coefficient action trivial.
A nonzero finite quotient of this homomorphism contradicts
F.24. Therefore finitely many points $x_1,\ldots,x_r$ suffice
to span $W$.

Subtract $\kappa(b_0)$ from $\xi$. The normalized class has
geometric functional
$\lambda:T_\ell A\to\Z_\ell(1)$ and is zero at $x_0$.
It factors through $W$ rationally. If $b_i$ is the actual value
at $x_i$, evaluation says
\[
 \lambda(c_{x_i}(h))=\kappa_\ell(b_i/b_0)(h).
\]
Choose one number field $M\supset F$ containing the coordinates
of every $x_i$, every $b_i$, and $b_0$, and put $L'=Mk$.
Restriction to $G_{L'}$ preserves the spanning property.
Both $W$ and $\Q_\ell(1)$ have trivial $G_{L'}$-action.
For $g\in G_M$ and $h\in G_{L'}$, the cocycle identity gives
\[
 c_{x_i}(ghg^{-1})=g c_{x_i}(h),\qquad
 \kappa_\ell(b_i/b_0)(ghg^{-1})
     =\chi_\ell(g)\kappa_\ell(b_i/b_0)(h).
\]
The points and actual values are $M$-rational, and $G_{L'}$
is normal in $G_M$, which are the hypotheses needed here.
Combining these formulas proves
$\lambda(gv)=\chi_\ell(g)\lambda(v)$ on each spanning image,
hence on all of $W$. Thus the original integral functional
$T_\ell A\to\Z_\ell(1)$ is $G_M$-equivariant.

Now apply the integral Tate-Hom theorem for semiabelian
varieties over the \emph{number field} $M$:
\[
 \operatorname{Hom}_M(A,\mathbb G_m)\otimes\Z_\ell
 \xrightarrow{\sim}
 \operatorname{Hom}_{G_M}(T_\ell A,\Z_\ell(1)).
\]
This is \citet[Lemma~6.3.2, pp.~81--82]{F_Jossen2009}, whose
proof and hypotheses have been opened. It is not being applied
directly over the infinite field $L$.
In the proper case the Hom-group is zero; equivalently, at a
good prime away from $\ell$, Frobenius eigenvalues on $T_\ell J$
have absolute value $q^{1/2}$ and those on $\Q_\ell(1)$ have
absolute value $q$, as in that source's Proposition~6.3.1.

The normalized unit lattice is independent of $M$ and $\ell$.
Indeed every geometric unit has divisor supported on the
$F$-rational boundary. Each Galois conjugate has the same divisor;
after normalization at $x_0$ it also has the same value there,
so the quotient is the constant one. Thus every normalized
geometric unit is already defined over $F$. This lattice is a
finitely generated free subgroup of the boundary-divisor lattice.
The integral Tate-Hom theorem places the normalized class in
its tensor product with $\Z_\ell$ for every $\ell$. Equality
of arithmetic classes follows from their equal geometric
functionals and zero values at $x_0$. One fixed basis therefore
gives the displayed coefficients in $\widehat\Z$.

There is no assumption that the lattice of principal boundary
divisors is primitive. The integral theorem gives the assertion
directly. Alternatively, after the rational Tate-Hom theorem,
the Kummer exact sequence gives saturation: the quotient of
integral $H^1$ by the completed unit group is
$T_\ell\operatorname{Pic}(U)$, which is torsion free. An
integral class in the rational unit span consequently lies in
the integral completed unit group. This explains why torsion
boundary classes do not create an unexamined denominator.
\end{proof}

\begin{claim}{F.27}{Proved}
Let $L/\Q^{\mathrm{solv}}$ be finite and let $U/L$ be a smooth
separated geometrically connected curve. Every compatible
multiplicative Kummer class with actual values at all points over
finite extensions is, after a finite constant extension, the class
of a regular unit $h$. At a point over any further finite field
$E$, its actual value divided by $h(x)$ belongs to
$(\Q^{\mathrm{solv}})^*$.
\end{claim}
\begin{proof}
Extend $L$ finitely to split the boundary and provide $x_0$.
Use the unit-lattice expression of F.26. For each boundary
point $d$, Riemann--Roch supplies a finite map
$\pi:C\to\mathbb P^1_L$ with a simple zero at $d$ and
nonzero values at the other boundary points. Delete the branch
values, $0,\infty$, and all boundary images. We obtain a finite
\emph{\'etale} map $W\to V$ with $W\subset U$ and $V$
a punctured rational line. The norm of the restricted class
has actual values: at a rational point over any finite extension,
its value is the product of the field norms of the given values
at the residual points of the fibre. F.18 makes this norm class
the class of an actual unit on $V$.

The residue at $0$ is the original residue $\lambda_d$.
The only boundary point over $0$ is the rational point $d$;
all interior points have residue zero. As in F.5, norms weight
residues by residue degree, not ramification index. Therefore
every $\lambda_d$ is an ordinary integer.
The divisor matrix of the basis units in F.26 is an injective
integer matrix. An invertible square minor and its adjugate
show that each profinite unit exponent is rational. Since
$\Q\cap\widehat\Z=\Z$, all exponents are ordinary integers.
The resulting unit, multiplied by the actual constant value at
$x_0$, represents the whole class. F.16 identifies the kernel
of its evaluations with $D(E^*)=(\Q^{\mathrm{solv}})^*$,
proving the last assertion.
\end{proof}

\begin{claim}{F.28}{Proved}
Let $L/\Q^{\mathrm{solv}}$ be finite and let $U/L$ be a smooth
separated geometrically connected curve. The geometric point
function of every pinned relative morphism
$U_{\mathrm{et}}\to(\mathbb G_{m,L})_{\mathrm{et}}$ is the
point function of one regular unit of $U$. If that function is
constant, the entire topos morphism is the algebraic constant
map.
\end{claim}
\begin{proof}
Put $k=\Q^{\mathrm{solv}}$ and write $q$ for the geometric
point function. Field-source reconstruction supplies actual
values. After a finite Galois constant extension, F.27 gives
$h\in\mathcal O(U)^*$ and
\[
 q(x)=c_x h(x),\qquad c_x\in k^*.
\]
We first work over this enlarged $L$.

Suppose $h=h_0$ is constant. Choose a nonzero algebraic $\beta$
with $\beta/h_0\notin k$ and extend constants to contain it.
The inverse image of $\beta$ is closed and has no closed point,
so it is empty. Pull back $t-\beta$ on all of $U$ and use
F.27 to obtain a unit $g$ with
\[
 q(x)-\beta=d_xg(x),\qquad d_x\in k^*.
\]
All necessary further constant extensions are finite. The left
side belongs to the current constant field for every geometric
$x$. If $g$ were nonconstant, a closed point of degree larger
than $\deg(g:C\to\mathbb P^1)$, available by F.24, would
have $g(x)$ outside that field, a contradiction. Thus $g$ is
constant. Comparing with one fixed $x_0$ gives
\[
 \frac{h_0c_x-\beta}{h_0c_{x_0}-\beta}\in k,
\]
which forces $c_x=c_{x_0}$ because $\beta/h_0\notin k$.
The point map is constant. Its factorization through that closed
subtopos and its relative structure give the full algebraic
constant map.

Now suppose $h$ is nonconstant. The point map cannot be
constant, by the same high-degree argument. For $b\in k^*$,
let $U_b$ be the nonempty open complement of $q^{-1}(b)$.
Apply F.27 to the pullback of $t-b$, obtaining, after a finite
constant extension if necessary,
\[
 q(x)-b=d_{x,b}h_b(x),\qquad
 d_{x,b}\in k^*,\quad h_b\in\mathcal O(U_b)^*.
\]
Each removed fibre is finite, since it is a proper closed
subset of a curve. For $b=1$, the rational functions
$c_xh-d_{x,1}h_1-1$ have one fixed bound on the degree of their
pole divisors. F.24 gives a closed point, outside the finite
exceptional set and over the current finite constant field,
whose degree exceeds this bound. Vanishing at that point forces
a rational identity, hence
\[
 h_1=Ah+B,\qquad A,B\in k^*.
\]
For any $x$ with $h(x)\notin k$, the relation
\[
 (c_x-d_{x,1}A)h(x)=1+d_{x,1}B
\]
forces $c_x=-A/B=:c$. Such $x$ cannot lie in the removed
fibre. Thus $q(x)=ch(x)$ whenever $h(x)\notin k$.

For arbitrary $b\in k^*$, choose a high-degree point also of
degree greater than $\deg h$, so its $h$-value is outside the
current constant field and hence outside $k$. The same bounded
pole-divisor argument yields
\[
 h_b=e_b(ch-b),\qquad e_b\in k^*.
\]
Because $h_b$ is a unit on $U_b$, every zero of $ch-b$ inside
$U$ belongs to $q^{-1}(b)$. At any remaining point with
$h(x)\in k^*$, choose $b=ch(x)$. This forces $q(x)=ch(x)$
there as well. The identities have coefficients in $k$ even
when a particular $U_b$ required further constants.

The reconstructed unit is unique because it agrees with the
point function on all geometric points. Galois equivariance of
the original point function therefore makes it invariant under
the initial finite Galois constant extension, so it descends
to $L$. The reconstructed unit also has the pulled-back Kummer
class: their difference vanishes at every finite-extension point,
so F.24 applied to each finite Kummer quotient makes it zero.
This proves the assertion without replacing pointwise
agreement by an unsupported full-topos comparison.
\end{proof}

\subsection{All-scheme reconstruction over the maximal solvable extension}

\begin{claim}{F.29}{Proved}
Let $L/\Q^{\mathrm{solv}}$ be finite and let $X/L$ be a smooth
geometrically connected quasi-projective variety. Every finite
quotient of $\pi_1(X)$ is detected surjectively by a closed-point
decomposition group.
\end{claim}
\begin{proof}
The zero-dimensional case is $X=\Spec L$, and the curve case
is F.24. Let $Y\to X$ be the connected finite \emph{\'etale}
Galois cover of a prescribed quotient, and suppose $\dim X=d\ge2$.
Embed $X$ locally closed in $\mathbb P^N_L$. Each geometric
component $Y_i$ of $Y$ is geometrically irreducible, since it
is smooth and connected, and its composite to $\mathbb P^N$
is quasi-finite with image $X_{\bar L}$. Its nonempty fibres
therefore all have dimension zero.
By \citet[Theorem~1.5]{F_PoonenSlavov2020}, the constructible
locus of hyperplanes whose pullback to $Y_i$ is not geometrically
irreducible has dimension at most
\[
 \operatorname{codim}(X_{\bar L},\mathbb P^N)+1
       =N-d+1<N.
\]
Remove the closures of these finitely many bad loci and the bad
smoothness locus for $X$. The remaining nonempty open is
Galois invariant and defined over $L$. Since $L$ is infinite,
it has an $L$-rational hyperplane. The section of $X$ is smooth
and geometrically connected, and every geometric component of
its pulled-back cover is still irreducible. The Galois group
continues to permute these components transitively, so the
pulled-back cover is connected over $L$. Thus arithmetic
constants, not only geometric monodromy, are retained.
Iterate to a smooth geometrically connected curve $C\subset X$
with $Y_C$ connected. F.24 gives a closed point of $C$ whose
decomposition group surjects onto the chosen Galois group.
As a point of $X$ it has the same residue field and supplies
the required detection.
\end{proof}

\begin{claim}{F.30}{Proved}
Let $L/\Q^{\mathrm{solv}}$ be finite and let $X/L$ be a smooth
geometrically connected affine variety with $x_0\in X(L)$.
Every compatible multiplicative Kummer class having actual
values at all finite-extension points is, after finite constants,
the class of an actual regular unit. Consequently the point
function of every pinned relative morphism
$X_{\mathrm{et}}\to(\mathbb G_{m,L})_{\mathrm{et}}$ is one
regular unit of $X$.
\end{claim}
\begin{proof}
The generalized Albanese $A_X$ pointed by $x_0$ has the
universal property
\[
 \operatorname{Hom}(A_X,\mathbb G_m)
 \simeq\{u\in\mathcal O(X)^*:u(x_0)=1\}
\]
\citep[Section~3]{F_SpiessSzamuely2003}. The exact cohomological
comparison needed here is integral. Over an algebraic closure,
that source's Proposition~4.1 gives
\[
 0\longrightarrow
 \operatorname{Hom}(A_X[\ell^r],\mu_{\ell^r})
 \longrightarrow H^1(X_{\bar L},\mu_{\ell^r})
 \longrightarrow \operatorname{Tor}(S,\Z/\ell^r)
 \longrightarrow0,
\]
where $S$ is the finite torsion subgroup of the N\'eron--Severi
quotient by boundary classes. The transition on its torsion
term is multiplication by $\ell$; its inverse limit is zero.
The finite left terms satisfy the Mittag--Leffler condition.
Passing to the limit therefore gives
\begin{equation}\label{eq:F30:albanese_integral}
 H^1(X_{\bar L},\Z_\ell(1))
 \simeq\operatorname{Hom}(T_\ell A_X,\Z_\ell(1)).
\end{equation}
Corollary~4.2 states this inverse-limit consequence explicitly.
Equivalently, Proposition~4.3 gives a finite kernel in
$\pi_1(X_{\bar L})^{\mathrm{ab}}\to T A_X$, killed by
homomorphisms to torsion-free $\Z_\ell(1)$.
All hypotheses hold: characteristic zero supplies a smooth
projective compactification by resolution, and every prime is
prime to the characteristic. The comparison is the one induced
by the pointed Albanese map, so evaluations are the corresponding
point Kummer cocycles.

Choose a smooth projective normal-crossings compactification
$\bar X$ preserving $X$, as in F.7. After a finite constant
extension its geometric boundary components are individually
defined. Descend $X,\bar X,x_0,A_X$ to a number field
$F\subset L$ with $L=F\Q^{\mathrm{solv}}$.
For each $\ell$, use the quotient
$(V_\ell A_X)_{G_L}$ and its integral image lattice.
The spanning, finite-number-field, and conjugation arguments
in F.26 apply verbatim, with F.29 replacing F.24 and
\eqref{eq:F30:albanese_integral} replacing the curve comparison.
More explicitly, a missing span would give a nonzero normalized
finite Kummer quotient invisible at all closed points, contradicting
F.29. Finitely many spanning points and their actual values
lie in one number field $M$. Restriction to $G_{M\Q^{\mathrm{solv}}}$
retains their spans, and conjugation by $G_M$ makes the geometric
functional arithmetic. Jossen's integral semiabelian Tate-Hom
theorem then expresses it in the normalized regular-unit lattice
\citep[Lemma~6.3.2]{F_Jossen2009}.

This is one fixed finitely generated free lattice: a normalized
geometric unit has divisor on the split boundary, and its Galois
conjugates have the same divisor and the same value one at $x_0$.
They are therefore equal. This supplies, for one basis
$u_1,\ldots,u_s$, an expression
\[
 \xi=\kappa(b_0)+\sum_j\lambda_j\kappa(u_j),
       \qquad \lambda_j\in\widehat\Z,
\]
with the actual value $b_0$ at $x_0$. This conclusion concerns
the class $\xi$, not vanishing of the whole Picard Tate module.

For a boundary component $D_i$, choose a smooth geometric point
on it away from the other components. After finite constants,
sufficiently ample complete intersections with specified first
jets give a smooth geometrically connected curve in $\bar X$
through that point and transverse to $D_i$. The finite-jet
Bertini predicates are precisely those checked in F.7, using
\citet[Lemmas~33.47.1, 33.47.3 and 37.32.3]{F_StacksBertini}.
Its open part in $X$ carries a point-theoretic restriction of
$\xi$. F.27 says this restriction is an actual unit class.
Its residue at the chosen boundary point is exactly the global
residue on $D_i$, because the intersection multiplicity is one.
At finite Kummer level this is the local calculation with a
parameter $z$ and, after \emph{\'etale} localization, a
representative $u z^a$; a transverse parameter pulls $z$ to a
unit times a parameter. Thus every global residue is an ordinary
integer. The injective integer divisor matrix of the $u_j$ and
$\Q\cap\widehat\Z=\Z$ force every $\lambda_j$ to be an
ordinary integer. We obtain an actual unit $h$ representing
$\xi$, normalized by $h(x_0)=b_0$.

Now let $\xi$ be the pullback of the coordinate class under the
pinned morphism in the statement, and let $q$ be its geometric
point function. Work first after the finite constant extension
just used. For any geometric point $x$, extend constants
finitely to make it rational. If $\dim X\ge2$, sufficiently
ample complete intersections through $x_0$ and $x$, with
independent first jets there, give a smooth geometrically
connected curve $C\subset X$. Smoothness and irreducibility
off the finite base locus and smoothness at the prescribed
points are the same finite-jet Bertini assertions used in F.7.
For $\dim X=1$ take $C=X$; dimension zero is the field-source
case. F.28 realizes the restricted point function by a regular
unit $v_C$ and gives
\[
 \kappa(v_C)=\xi|_C=\kappa(h|_C).
\]
The unit $v_C/h|_C$ has zero compatible Kummer class. Its
ordinary divisor is divisible by every integer and hence zero,
so it is constant; F.16 puts that constant in
$(\Q^{\mathrm{solv}})^*$. Its value at $x_0$ is one, so
$q(x)=h(x)$. This also works when $h|_C$ is constant.
Thus the entire point function equals $h$.
Finally uniqueness on geometric points and Galois equivariance
descend $h$ through the initial finite constant extension.
\end{proof}

\begin{claim}{F.31}{Proved}
Every finite extension $L/\Q^{\mathrm{solv}}$ has
closed-point-preserving \emph{\'etale} reconstruction. Precisely,
for all finite-type $L$-schemes $X,Y$, the natural map
\[
 \operatorname{Hom}_L(X^{\mathrm{awn}},Y)
 \xrightarrow{\sim}
 \pi_0\operatorname{Hom}^{\mathrm{pin}}_{BG_L}
       (X_{\mathrm{et}},Y_{\mathrm{et}})
\]
is bijective, with pinned understood in its actual
closed-point-preserving sense. This supplies the all-scheme
property for $\Q^{\mathrm{solv}}$ itself. It does not classify
all characteristic-zero reconstruction fields intrinsically.
\end{claim}
\begin{proof}
Apply \citet[Theorem~4.18, equivalence of (1) and (5)]{F_CHW2024}.
Its condition (5) asks that, for every connected regular affine
finite-type source over $L$ having an $L$-point, every pinned
relative morphism to $\mathbb G_m$ have the point function of
one regular unit. In characteristic zero, regular finite-type
sources over a field are smooth. Connectedness together with
an $L$-point makes such a smooth source geometrically connected:
its geometric connected components form a transitive Galois set,
and the component containing that rational point is fixed.
F.30 supplies exactly the required unit on exactly these sources.
The theorem then reconstructs the full small-topos Hom-groupoids
using its affine-coordinate, \emph{\'etale}-chart, and $h$-descent
arguments; no full-topos conclusion is inferred merely from
pointwise agreement. Characteristic-zero absolute weak
normalization is seminormalization, as specified in F.7.
The relative-only interpretation of R27 remains refuted by
F.1 and F.17. The new theorem proves its corrected version
over this additional field class while leaving F.9's intrinsic
classification request unresolved.
\end{proof}


\subsection{Algebraic descent and finitely generated extensions of the solvable closure}

\begin{claim}{F.32}{Proved}
Let $K/k$ be an algebraic extension of characteristic-zero fields.
If $K$ satisfies closed-point-preserving \emph{\'etale}
reconstruction, then so does $k$. The extension need not be
finite or Galois. In particular F.31 also gives reconstruction
over every subfield of a finite extension of $\Q^{\mathrm{solv}}$.
\end{claim}
\begin{proof}
We first construct the scalar extension of a relative geometric
morphism $\phi:X_{\mathrm{et}}\to Y_{\mathrm{et}}$, where $X,Y$
are finite-type $k$-schemes. For each finite subextension
$E/k$ of $K/k$, the relative structure identifies
$\phi^*(Y_E)$ with $X_E$. Slicing at this finite \emph{\'etale}
object gives
\[
 \phi_E:(X_E)_{\mathrm{et}}\longrightarrow(Y_E)_{\mathrm{et}}
 \quad\text{over }BG_E.
\]
Successive slices give canonical compatibilities for $E\subset E'$.
Let $\mathcal C(Z)$ denote the category of quasicompact separated
\emph{\'etale} $Z$-schemes. By \citet[Lemma~A.4,
Theorem~A.8 and Corollary~A.9]{F_CHW2024}, this is a basis
and $\phi_E^*$ maps $\mathcal C(Y_E)$ into $\mathcal C(X_E)$.
This assertion is stronger than preservation of coherent sheaves.
The proof of Theorem~A.8 uses the complementary clopen $V'$ to
the diagonal of a separated chart $V$: the decomposition
$V\amalg V'\simeq V\times V$ pulls back to
$W\amalg W'\simeq W\times W$. The representing constructible
\emph{\'etale} algebraic space $W$ therefore has closed diagonal.
It is separated and locally quasi-finite, hence a scheme;
coherence supplies quasicompactness. Thus nonseparated
constructible objects cause no change to the stated basis.

Finite-presentation descent gives
\[
 \mathcal C(X_K)\simeq\mathop{\mathrm{2\!\!\ colim}}_E
       \mathcal C(X_E),\qquad
 \mathcal C(Y_K)\simeq\mathop{\mathrm{2\!\!\ colim}}_E
       \mathcal C(Y_E).
\]
The exact affine-site statement is Stacks, Lemma~59.51.2,
Tag~0EYL; Lemma~32.10.1 descends objects, morphisms and their
finite relations \citep{F_StacksLimits}. The \emph{\'etale}
and separated predicates descend after increasing $E$.
Finite coverings and their surjectivity also descend.
The same conclusions for the displayed basis follow by finite
affine covers and their overlap data. Successive-slice
compatibility gives
\[
 \phi_{E'}^*(V\times_E E')\simeq\phi_E^*(V)\times_E E'.
\]
Consequently the finite-stage inverse images define a functor
$\mathcal C(Y_K)\to\mathcal C(X_K)$ preserving finite limits
and covering families: every such finite diagram is represented
at one finite stage. Stacks, Proposition~7.14.7, Tag~00X6,
therefore gives a geometric morphism $\phi_K$ on the full
small topoi. Its structure over $BG_K$ is inherited on finite
\emph{\'etale} constant objects, each of which descends to a
finite stage. This construction commutes with further algebraic
extension and with the projection to $\phi$. It does not assert
an arbitrary fibre-product theorem for small \emph{\'etale} topoi.

If $\phi$ is pinned, so is $\phi_K$. A closed $x_K\in X_K$
projects to a closed $x\in X$, because $\kappa(x)$ is algebraic
and finitely generated over $k$. The image of $x_K$ under
$\phi_K$ lies over the closed point $\phi(x)$.
That fibre is $\Spec(\kappa(\phi(x))\otimes_kK)$, finite over
$K$, so its points are closed. For an algebraic closure
$\bar k\supset K$, field-source reconstruction
\citep[Proposition~3.6]{F_CHW2024} moreover gives
\[
 \phi_K(\bar k)(x)=\phi(\bar k)(x)
 \qquad(x\in X(\bar k)).
\]
The retained relative structure is used in this equality.

Now take $X=\Spec A$ regular connected affine with a $k$-point
and $Y=\mathbb G_m$. The source is smooth and geometrically
connected, so $X_K$ has the same properties and a $K$-point.
Reconstruction over $K$ gives a unit
$h_K\in(A\otimes_kK)^*$ with the point function of $\phi_K$.
Its image $\bar h\in A\otimes_k\bar k$ has the
$G_k$-equivariant point function of $\phi$. Thus
$\sigma\bar h$ and $\bar h$ agree at every $\bar k$-point,
for each $\sigma\in G_k$. Reducedness gives
$\sigma\bar h=\bar h$. The invariant ring
$(A\otimes_k\bar k)^{G_k}=A$: express an element using finitely
many $k$-linearly independent elements of $A$, and invariance
fixes each coefficient. The inverse of $\bar h$ is invariant
as well. Hence $h_K$ descends to a unit of $A$ with precisely
the original point function. CHW, Theorem~4.18(5), now gives
full reconstruction over $k$. No descent of a chosen topos
isomorphism or of a merely pointwise comparison is assumed.
Even when $K/k$ is not normal, $\sigma\bar h$ is well defined
over $\bar k$, so the proof is unchanged.
\end{proof}

\begin{claim}{F.33}{Proved}
Put $k=\Q^{\mathrm{solv}}$. For every finitely generated
extension $L/k$,
\[
 D(L^*)=\bigcap_{n\ge1}(L^*)^n=k^*.
\]
There is a finite Galois number field $N/\Q$ not contained in
$L$. Any finite list of schemes, maps, points and coefficients
over $L$ can be descended to a finitely generated field
$F/\Q$ with $L=Fk$; then $L/F$ is Galois and prosolvable.
\end{claim}
\begin{proof}
The algebraic closure $k'$ of $k$ in $L$ is finite over $k$
\citep[Lemma~9.26.11]{F_StacksFields}. Choose a normal projective
model $Z$ of $L/k$. Normalization is finite in this
characteristic-zero finite-type setting. The global regular
functions of $Z$ form $k'$: properness makes them algebraic,
and every element of $k'$ is integral in every normal local ring.
If $a\in D(L^*)$, its valuation at each prime divisor of $Z$
is divisible by every positive integer, hence is zero. Normality
makes $a$ and $a^{-1}$ global regular functions. Any root
$z\in L$ with $z^n=a$ also has zero divisor, so lies in $k'$.
Thus $a\in D(k'^*)=k^*$ by F.16. The reverse inclusion follows
from divisibility of $k^*$. This also proves the assertion over
every finite extension of $L$.

If $\bar\Q\subset L$, then $\bar\Q\subset k'$, contrary to
F.18's unbounded degrees of finite extensions of $k$.
Choose an algebraic number outside $L$ and take its finite
Galois closure $N/\Q$. To descend the finite data, start with
finitely many generators of $L/k$ and adjoin all the finitely
many coefficients and algebraic constants occurring in the data
and their relations. The resulting $F\subset L$ is finitely
generated over $\Q$ and satisfies $L=Fk$.
The extension $Fk/F$ is the scalar extension of the Galois
prosolvable extension $k/\Q$. Its Galois group injects as a
closed subgroup of $\operatorname{Gal}(k/\Q)$ by restriction,
so is prosolvable. In particular every finite Galois subextension
of $L/F$ is solvable.
\end{proof}

\begin{claim}{F.34}{Proved}
Let $L/\Q^{\mathrm{solv}}$ be finitely generated. Every finite
arithmetic quotient of the fundamental group of a smooth
separated geometrically connected curve $U/L$ is detected
surjectively by a closed-point decomposition group. Such points
can be chosen with unbounded degrees. Every finite quotient
on a smooth geometrically connected quasiprojective variety over
$L$ is likewise detected by a closed point.
\end{claim}
\begin{proof}
We specify the arithmetic changes to F.24, whose geometric
pencil and finite-group calculations apply unchanged.
Descend the curve, completion, connected finite $G$-torsor,
its full $G$-action and torsor identity to a finitely generated
$F/\Q$ with $L=F\Q^{\mathrm{solv}}$. Let $M/F$ be the finite
Galois constant field of the descended cover, and write
$1\to N\to G\to P\to1$, with $P=\operatorname{Gal}(M/F)$.
Connectedness after extension to $L$ implies $M\cap L=F$.
The length-at-most-four jet argument constructs the same
simple-branch degree-$n$ pencil over $F$ for unbounded $n$.
The field $F$ embeds in $\C$, so the SGA~1 comparison and the
branch-loop calculation give the same group
\[
 H\subseteq(G\times_P\cdots\times_PG)\rtimes S_n.
\]
In particular $A_n\subset H^\infty$ and
$H^\infty\cap N^n$ projects onto $N$ at every coordinate.
These are the proved spanning-tree and support-two-commutator
calculations of F.24; they neither use that $F$ is a number
field nor assert that the whole wreath product occurs.
Possible extra constants in the Galois closure remain harmless.

The field $F$ is Hilbertian. Here is the deduction from the
opened source. Hilbert's theorem over $\Q$ holds in every
affine dimension \citep[Theorem~3.4.1]{F_Serre2007}.
For a finite family of irreducible polynomials over
$\Q(t_1,\ldots,t_d)(T)$, make them monic in the root variable
and clear the coefficient denominators. Simultaneous Hilbert
specialization of $(t_1,\ldots,t_d,T)$ over $\Q$, avoiding
those denominators, gives irreducible specializations.
The chosen value of $T$ therefore leaves each polynomial
irreducible over $\Q(t_1,\ldots,t_d)$: a monic factorization
over that fraction field has coefficients in the integrally
closed localized polynomial ring, and would specialize at the
chosen $t$-value to a factorization of the same positive degrees.
Thus the rational-function field is Hilbertian. A finitely
generated $F/\Q$ is finite over such a field, and finite
extensions preserve Hilbertianity
\citep[Corollary~3.2.2]{F_Serre2007}.

Proposition~3.3.1 of the same source now specializes the
Galois closure with its full group $H$ over $F$. For the
specialized extension $\Omega/F$, set
$J=\operatorname{Gal}(\Omega L/L)\subset H$.
F.33 gives $H^\infty\subset J$, and $M\cap L=F$ gives
$J\twoheadrightarrow P$. The coordinate stabilizer
$J_i=J\cap\operatorname{Stab}(i)$ projects onto $G$:
its image contains $N$ by the commutator calculation, and
literal $A_n$ moves the block image of an element of $J$
back to $i$ without changing its prescribed $P$-image.
Transitivity of $A_n$ makes the degree-$n$ fibre one closed
point, whose arithmetic decomposition group has image $J_i$.
No perfectness of $P$ or $J$ is assumed.

For a higher-dimensional source use F.29's hyperplane
construction on every geometric component of the given finite
connected cover. The Poonen--Slavov theorem applies because
each composite with the projective embedding is quasi-finite;
the source dimension is at least two until a curve is reached.
The simultaneous good opens have rational points over the
infinite field $L$. Geometric components remain irreducible
and their arithmetic permutation remains transitive, so the
curve retains the entire finite quotient. The curve case
then supplies the detecting point.
\end{proof}

\begin{claim}{F.35}{Proved}
Let $k=\Q^{\mathrm{solv}}$, let $L/k$ be finitely generated,
and let $U\subset\mathbb G_{m,L}$ be nonempty open. Suppose
$\xi\in\varprojlim_n H^1(U,\mu_n)$ has actual multiplicative
Kummer values at every point over every finite extension $E/L$.
Then $\xi=\kappa(h)$ for an actual regular unit
$h\in\mathcal O(U)^*$, and every prescribed value $b_x$ satisfies
$b_x/h(x)\in k^*$.
\end{claim}
\begin{proof}
Write $P_0=T$ and let $P_1,\ldots,P_r$ be the other monic
irreducible puncture polynomials. The coordinate ring is a
localization of a principal ideal domain. Choose $x_0\in U(L)$
and an actual value $b_0$ there. The Kummer sequence gives
\[
 \xi=\kappa(b_0)+\sum_{j=0}^r\lambda_j
       \bigl([P_j(T)]-[P_j(x_0)]\bigr),
       \qquad\lambda_j\in\widehat\Z.
\]
Take $F_0/\Q$ finitely generated inside $L$ with $L=F_0k$,
containing the coefficients, $x_0,b_0$, and fix a transcendence
basis $t_1,\ldots,t_d$ of $F_0/\Q$.
Choose a finite Galois number field $N\not\subset L$ by F.33.
Take a finite Galois extension $M/\Q(t_1,\ldots,t_d)$ containing
$F_0,N$ and all roots of every $P_j$. The extension $Mk/L$
is finite Galois and has the nontrivial Galois subextension
$LN/L$. Choose $\sigma\in\operatorname{Gal}(Mk/L)$ acting
nontrivially on $N$. Normality of $M/\Q(t_1,\ldots,t_d)$
makes $M$ stable under $\sigma$. Put $C=M\cap\bar\Q$.
This is a finite Galois number field: embeddings of $C/\Q$
extend while fixing the independent $t$-variables, and normality
preserves $M$. Moreover $\sigma|_C\ne1$, since $N\subset C$.
The polynomials' coefficients lie in $F_0$ and are fixed by
$\sigma$; individual roots need not be fixed.

For $d=0$ use F.18's completely split primes. For $d>0$ the
following construction supplies the required valuations.
Choose a root $r_i$ of each $P_i$. There is a normal finite
\emph{\'etale} model $B$ of $M/C(t_1,\ldots,t_d)$ over
\[
 A=\mathcal O_C[S^{-1}][t_1,\ldots,t_d,1/f],
 \qquad0\ne f\in\Z[t_1,\ldots,t_d],
\]
where $S$ is a finite set of rational primes. Indeed choose a
primitive element, clear its polynomial denominators, and invert
its nonzero discriminant. The resulting monogenic algebra is
finite \emph{\'etale} and normal. Inverting further norms
makes all the finitely many required elements integral or units.
Products of conjugate denominators and enlargement of $S$
allow $f$ to have integral rational coefficients. Thus $B$ is
the integral closure of $A$ in $M$ and is $\sigma$-stable.
Arrange that all coefficients and roots are integral and that
\[
 P_i'(r_i),\quad P_j(r_i)\ (j\ne i)
\]
are units. These elements are nonzero by separability and
distinctness of the puncture polynomials.

Take pairwise distinct rational primes $p_i$ splitting
completely in $C$, outside $S$, with $f\bmod p_i\ne0$.
F.18 supplies infinitely many such split primes without a
Chebotarev hypothesis. For a prime $\mathfrak p_i$ of $C$
over $p_i$, the localization $A_{\mathfrak p_iA}$ is the
Gauss discrete valuation ring with uniformizer $p_i$ and
residue field $\F_{p_i}(t_1,\ldots,t_d)$. A height-one prime
of $B$ above it gives an unramified prolongation $v_i$ to $M$,
so $v_i(p_i)=1$ and its residue field is infinite.
Complete splitting and $\sigma|_C\ne1$ imply
$v_i\ne\sigma^{-1}v_i$, since their restrictions to $C$ are
different primes. All pairs belonging to different $i$ are
disjoint because their rational residue characteristics differ.

Localizing $B$ at the complement of these finitely many
height-one primes gives a normal noetherian one-dimensional
semilocal ring $R$, whose local rings are the selected discrete
valuation rings. The Chinese remainder theorem supplies
$\alpha\in R$ such that
\[
 \alpha\equiv r_i+p_i\pmod{\mathfrak p_{v_i}^2},\qquad
 \alpha\equiv c_i\pmod{\mathfrak p_{\sigma^{-1}v_i}},
\]
where $c_i$ avoids every root of every reduced $P_j$.
The residue fields are infinite for $d>0$; for $d=0$ take
$p_i>\sum_j\deg P_j$, as in F.18. Taylor expansion and the
unit conditions give
\[
 v_i(P_j(\alpha))=\delta_{ij},\qquad
 v_i(P_j(\sigma\alpha))
       =v_{\sigma^{-1}v_i}(P_j(\alpha))=0.
\]
Hence $\alpha\in U(Mk)$ and
$R_j=P_j(\sigma\alpha)/P_j(\alpha)$ has
$v_i(R_j)=-\delta_{ij}$.

Choose one finite number field $S_0\subset k$ such that
$F=MS_0$ contains the actual value $b_\alpha\in Mk$.
This $F$ is fixed for all radical levels and is $\sigma$-stable.
If $z^n=a\in F^*$ with $z\in Mk$, every
$\tau\in\operatorname{Gal}(Mk/F)$ satisfies
$\tau z/z\in\mu_n\subset k$. The automorphisms $\tau$ and
$\sigma$ commute: check on $M$, which $\tau$ fixes and
$\sigma$ preserves, and on $k$, which $\sigma$ fixes and
$\tau$ preserves. Therefore $\sigma z/z$ is fixed by every
such $\tau$. Since $Mk/F$ is Galois, its fixed field is $F$,
and
\[
 \sigma a/a\in(F^*)^n.
\]
This radical descent uses a finitely generated $F/\Q$, not a
number-field hypothesis.

For integer representatives $a_{j,n}$ of $\lambda_j\bmod n$,
the evaluation at $\alpha$ and this descent give
\[
 \frac{\sigma b_\alpha}{b_\alpha}
       \prod_jR_j^{-a_{j,n}}\in(F^*)^n.
\]
The integer characters $h_i=v_i\circ N_{F/M}$ satisfy
$h_i(R_j)=-[F:M]\delta_{ij}$, so
\[
 [F:M]\lambda_i=-h_i(\sigma b_\alpha/b_\alpha)\in\Z
 \quad\text{inside }\widehat\Z.
\]
Thus $\lambda_i\in\Q\cap\widehat\Z=\Z$. The desired unit is
\[
 h(T)=b_0\prod_j
       \left(\frac{P_j(T)}{P_j(x_0)}\right)^{\lambda_j}.
\]
The residual ratio belongs to $D(E^*)=k^*$ by F.33.
\end{proof}

\begin{claim}{F.36}{Proved}
Let $L/\Q^{\mathrm{solv}}$ be finitely generated and let $X/L$
be smooth, geometrically connected and affine with $x_0\in X(L)$.
A compatible Kummer class on $X$ having actual multiplicative
values at every finite-extension point is the class of an actual
regular unit, normalized by its value at $x_0$.
Every pinned relative geometric morphism
$X_{\mathrm{et}}\to(\mathbb G_{m,L})_{\mathrm{et}}$ has the
geometric-point function of a unique regular unit on $X$.
\end{claim}
\begin{proof}
We give the arithmetic replacements in F.26--F.30 explicitly.
After a finite extension splitting the boundary of a smooth
compactification, descend $X,x_0$, its generalized Albanese $A$,
and the needed geometric data to a finitely generated $F/\Q$
with $L=F\Q^{\mathrm{solv}}$. For curves use the generalized
Jacobian. Spie\ss--Szamuely's integral comparison used in F.30
identifies
\[
 H^1(X_{\bar L},\Z_\ell(1))
   \simeq\operatorname{Hom}(T_\ell A,\Z_\ell(1)).
\]
Its hypotheses are smooth quasiprojective geometry and a smooth
projective compactification, available in characteristic zero;
no number-field hypothesis occurs. The finite kernel in geometric
abelian $\pi_1$ is killed by the torsion-free target.

Put $H=G_L\triangleleft G_F$ and
$W=(V_\ell A)_H$, with lattice the image of $T_\ell A$.
The normalized point Kummer homomorphisms into this lattice
span $W$. Otherwise an annihilating integral functional would
produce a nonzero finite quotient of $\pi_1(X)$ trivial on
every closed-point decomposition group, contradicting F.34.
Here the rational point splits the arithmetic group and
$\mu_\infty\subset L$ makes the normalized functional an
actual arithmetic homomorphism. Choose finitely many point
images spanning $W$. Their coordinates and actual multiplicative
values lie in one finite extension $M/F$, after adjoining the
finitely many constants from $\Q^{\mathrm{solv}}$ occurring
in their equations. Restriction to
$G_{M\Q^{\mathrm{solv}}}$ preserves their rational spans,
because it has finite index in their original point groups and
the projected Kummer maps are homomorphisms.

The conjugation identities of F.26 now show that the functional
$\lambda:T_\ell A\to\Z_\ell(1)$ is $G_M$-equivariant.
Indeed, for each selected point $i$, all
$g\in G_M$ and $h\in G_{M\Q^{\mathrm{solv}}}$ satisfy
\[
 c_i(ghg^{-1})=g c_i(h),\qquad
 \kappa(\beta_i/b_0)(ghg^{-1})
       =\chi(g)\kappa(\beta_i/b_0)(h).
\]
The equality $\lambda c_i=\kappa(\beta_i/b_0)$ gives the
required equivariance on their spanning images and hence on
$W$. Normality of $G_L$ in $G_F$ makes this quotient
$G_M$-stable. Thus the functional is equivariant over the
finitely generated field $M$, before any Tate-Hom theorem is used.

The necessary primary theorem over this field is
\citet[Theorem~4.3(ii), p.~676]{F_Jannsen1995}:
for 1-motives over a finitely generated field, the natural map
\[
 \operatorname{Hom}_M(M_1,M_2)\otimes\Z_\ell
 \xrightarrow{\sim}
 \operatorname{Hom}_{G_M}(T_\ell M_1,T_\ell M_2)
\]
is an isomorphism. The original proof reduces by extension
sequences to the abelian Tate theorem and the Mordell--Weil
injection for the dual abelian variety. The same statement is
recalled in \citet[\S1.4, p.~6]{F_Jossen2014}; its reference
to Jannsen's ``Theorem~4.6'' differs from the opened original's
numbering~4.3(ii). Apply it to $[0\to A]$ and
$[0\to\mathbb G_m]$.
After boundary components are split and units are normalized
at $x_0$, the geometric unit lattice is fixed over $F$:
its divisors are fixed, and value~1 removes the scalar ambiguity.
The Albanese universal property therefore gives
\[
 \xi=\kappa(b_0)+\sum_j\lambda_j\kappa(u_j),
       \qquad\lambda_j\in\widehat\Z,
\]
in one normalized integral unit basis. This kills the Picard
Tate component of this point-theoretic class, not the whole
Picard Tate module.

For a curve, the Riemann--Roch map used in F.27, with a simple
zero at one rational boundary point and no other boundary zero,
reduces each residue to the norm of a class on a punctured line.
Norms preserve actual point values, and F.35 makes the residue
an ordinary integer. The comparison uses residue degrees,
not ramification indices, and the chosen boundary point
contributes degree one. The injective integral divisor matrix
of the unit basis then forces every $\lambda_j$ to be an
ordinary integer.
For higher dimension, choose a smooth transverse complete-
intersection curve at a smooth point of each geometric boundary
divisor, away from the other components, using the finite-jet
Bertini construction in F.7. Its actual-unit class has the same
boundary residue: at each finite level the local class is
$u z^a$ after \emph{\'etale} localization, and transversality
pulls $z$ back to a uniformizer. Thus all global residues are
integral, and the same divisor matrix gives the actual unit $h$.
Normalize it by $h(x_0)=b_0$.

The finite extensions used here do not create a descent gap.
Over a finite Galois extension making the construction, each
conjugate of $h$ has the same Kummer class. Their quotient has
zero divisor on a proper completion, hence is constant; its
value at $x_0$ is~1. Thus $h$ descends to $L$.
At each finite Kummer level the remaining difference from $\xi$
is geometrically zero, so is a constant class by the low-degree
Hochschild--Serre sequence. Evaluation at $x_0$ kills that
class and splits the constant-class inclusion. Hence
$\xi=\kappa(h)$ already over $L$.

For a pinned map on a curve, F.33 gives
$\phi(x)=c_x h(x)$ with $c_x\in(\Q^{\mathrm{solv}})^*$.
The point-function argument of F.28 now has all its hypotheses:
F.34 supplies points of unbounded degrees, and F.35 together
with the curve calculation supplies actual units on every
punctured source. If $h$ is constant, pass to a finite extension
containing $\beta$ with $\beta/h\notin\Q^{\mathrm{solv}}$.
The point values avoid $\beta$, so $T-\beta$ pulls back on
the whole source. The corresponding unit must be constant,
since a nonconstant rational function cannot take values in the
base field at closed points of degrees exceeding its degree.
Comparison with the value at $x_0$ makes the point function
constant. If $h$ is nonconstant, the actual unit from $T-1$
and a point of degree greater than the fixed pole bound give
$h_1=A h+B$, with $A,B\in(\Q^{\mathrm{solv}})^*$.
High-degree values fix the scalar $c=-A/B$; applying the same
argument to $T-b$ for each $b\in(\Q^{\mathrm{solv}})^*$ gives
$h_b=e_b(ch-b)$. Its zeros force $\phi(x)=ch(x)$ also where
$h(x)\in\Q^{\mathrm{solv}}$. The needed opens are nonempty
because the point function is nonconstant. The closed-subtopos
argument in the constant case and the normalization at $x_0$
give the unique unit. These arguments are the explicit
pole-degree and coefficient comparisons proved in F.28;
only their field predicates have changed. F.34 also shows
that its Kummer class is the original pulled-back class.

Finally, for any geometric point $x$ on $X$, make it rational
by a finite extension and take a smooth geometrically connected
curve through $x$ and $x_0$. The curve's point-function unit
and $h$ restricted to it have the same Kummer class. Their
ratio has zero divisor on the proper completion, hence is
constant, and its value at $x_0$ is~1. This proves
$\phi(x)=h(x)$ for every $x$. The unique unit descends through
all the finite extensions; dimension zero is field-source
reconstruction and dimension one was treated above.
\end{proof}

\begin{claim}{F.37}{Proved}
Let $k_0$ be any subfield of a finitely generated extension of
$\Q^{\mathrm{solv}}$. Then for all finite-type $k_0$-schemes
$X,Y$, the canonical map
\[
 \operatorname{Hom}_{k_0}(X^{\mathrm{awn}},Y)
 \longrightarrow
 \pi_0\operatorname{Hom}^{\mathrm{cp}}_{BG_{k_0}}
       (X_{\mathrm{et}},Y_{\mathrm{et}})
\]
is bijective. This sufficient field class contains the corrected
LM5-F.50 class. It does not supply the intrinsic necessary-and-
sufficient classification requested in F.9.
\end{claim}
\begin{proof}
First let $L/\Q^{\mathrm{solv}}$ be finitely generated.
For a connected regular affine finite-type $L$-scheme with an
$L$-point, characteristic zero gives smoothness, and connectedness
together with that point gives geometric connectedness.
F.36 supplies precisely the regular-unit point-function condition
in \citet[Theorem~4.18(5)]{F_CHW2024}. That theorem supplies
full \emph{\'etale}-chart reconstruction and $h$-descent,
hence all-scheme reconstruction over $L$.

Now let $k_0\subset L$ for such an $L$ and put
$K=k_0\Q^{\mathrm{solv}}\subset L$.
This intermediate field is finitely generated over
$\Q^{\mathrm{solv}}$. To see this, choose a finite
transcendence basis $t$ of $K/\Q^{\mathrm{solv}}$.
Then $K$ is algebraic over $\Q^{\mathrm{solv}}(t)$ inside the
finitely generated field extension $L/\Q^{\mathrm{solv}}(t)$.
The relative algebraic closure is finite by
\citet[Lemma~9.26.11]{F_StacksFields}, so $K$ is finite over
$\Q^{\mathrm{solv}}(t)$. The first paragraph gives
reconstruction over $K$. Since $K/k_0$ is algebraic, F.32
descends it to $k_0$.

A finitely generated extension of $\Q^{\mathrm{ab}}$ embeds
in its compositum with $\Q^{\mathrm{solv}}$ in an algebraic
closure; that compositum is finitely generated over
$\Q^{\mathrm{solv}}$. Thus the proved class contains F.50.
All pinning here includes closed-point preservation.
The literal relative-only definition remains refuted by F.1
and F.17. No implication characterizing all reconstruction
fields from the present sufficient class has been proved.
\end{proof}


\subsection{Intrinsic necessary conditions and excluded field families}

\begin{claim}{F.38}{Proved}
If $E/k$ is a finite extension of characteristic-zero fields,
then $k$ satisfies closed-point-preserving \emph{\'etale}
reconstruction if and only if $E$ does.
\end{claim}
\begin{proof}
Descent is F.32. Suppose conversely that $k$ reconstructs.
Finite-type $E$-schemes $X,Y$ are finite-type $k$-schemes.
A pinned morphism $\phi:X_{\mathrm{et}}\to Y_{\mathrm{et}}$
over $BG_E$ is also pinned over $BG_k$. Reconstruction over
$k$ gives a unique $k$-map $f:X^{\mathrm{awn}}\to Y$ and
an isomorphism to $\phi$ over $BG_k$.
The composite with $Y\to\Spec E$ and the given structure
$X^{\mathrm{awn}}\to\Spec E$ induce isomorphic morphisms
of topoi over $BG_k$. Faithfulness over $k$ makes these
algebraic maps equal. Hence $f$ is $E$-linear.

The comparison isomorphism is also over $BG_E$.
Indeed $BG_E$ is the slice of $BG_k$ at the finite
\emph{\'etale} object $\Spec E$; the structural maps from
$X$ and $Y$ are its distinguished sections after pullback.
Naturality of the $BG_k$-relative isomorphism respects those
sections. Equivalently, the groupoid of morphisms to this
finite \emph{\'etale} slice is discrete, so the two structural
comparisons cannot differ by a further $2$-cell.
For the latter assertion one can apply
\citet[Proposition~2.22]{F_CHW2024} to the finite-type target
$\Spec E$ over $k$; this is its actual stated scope.
Absolute weak normalization depends on the underlying scheme,
not on whether its structure is viewed over $k$ or $E$.
This proves fullness and faithfulness over $E$.
No ascent through an infinite algebraic extension is claimed.
\end{proof}

\begin{claim}{F.39}{Proved}
If a characteristic-zero field $k$ satisfies reconstruction,
then its absolute Galois group is slim: every open subgroup
has trivial centre.
\end{claim}
\begin{proof}
Fix $\bar k$ and $z\in Z(G_k)$, and put $X=\mathbb G_{m,k}$.
For every continuous discrete $G_k$-set $S$, the map
$s\mapsto zs$ is $G_k$-equivariant and natural because $z$
is central. It therefore gives a natural automorphism of the
structural pullback $p^*:BG_k\to X_{\mathrm{et}}$.
Equip the underlying identity geometric morphism of
$X_{\mathrm{et}}$ with this new structural $2$-cell, or its
inverse according to the chosen direction of that cell.
Call the resulting relative morphism $\theta_z$.
Its underlying topological map is the identity, so it is pinned.

Choose the inverse convention so that its relative geometric-point
function is
\[
 \theta_z(\bar k):\bar k^*\longrightarrow\bar k^*,
             \qquad a\longmapsto z(a).
\]
Here is the verification with the relative isomorphism classes
retained. Place $a$ in a finite Galois extension $E/k$.
The constant object is the $G_k$-set
$\operatorname{Hom}_k(E,\bar k)$, and the structural twist
acts on its embeddings by $\iota\mapsto z\circ\iota$.
Slicing the point morphism at this object changes exactly its
residue embedding by $z$. By
\citet[Proposition~3.6]{F_CHW2024}, relative pinned morphisms
$\Spec(\bar k)_{\mathrm{et}}\to X_{\mathrm{et}}$ are precisely
$X(\bar k)$, so the changed relative point is $z(a)$.
In particular the two relative point morphisms are isomorphic
only if $z(a)=a$.

Every $k$-point is fixed. An algebraic $k$-map $X\to X$
has the form $T\mapsto cT^m$, $c\in k^*$, $m\in\Z$.
The fixed values at $1$ and $2$ force $c=1$ and $m=1$.
Such a map fixes every geometric point. If $z\ne1$, an
$a\in\bar k^*$ with $z(a)\ne a$ gives a contradiction.
A global natural automorphism of the identity topos cannot
absorb the twist in the relative $2$-isomorphism quotient:
its restriction would identify the distinct relative point
morphisms just separated by Proposition~3.6. No computation
of the automorphisms of the entire topos is assumed.
Thus $Z(G_k)=1$. F.38 applies the same argument to every
finite extension $E/k$; its $G_E$ is the corresponding open
subgroup of $G_k$. This is slimness.
\end{proof}

\begin{claim}{F.40}{Proved}
For every characteristic-zero reconstruction field $k$ and every
finite extension $E/k$, the natural map
\[
 \operatorname{Aut}(E)\longrightarrow\operatorname{Out}(G_E)
\]
is injective. Together with F.39, this is equivalent to the
following intrinsic necessary condition, for a fixed algebraic
closure $\bar k$:
\[
 C_{\operatorname{Aut}(\bar k)}(G_E)=\{1\}
             \qquad\text{for every finite }E/k.
\]
No sufficiency assertion is made.
\end{claim}
\begin{proof}
It suffices by F.38 to treat $E=k$.
Extending a field automorphism $\sigma$ to $\bar k$ gives an
automorphism of $G_k$ by conjugation. Different extensions
differ by $G_k$, so its outer class is well defined.
Suppose $\sigma\ne1$ has trivial outer action. Inner-conjugate
group homomorphisms induce naturally isomorphic geometric
morphisms on classifying topoi: on continuous Galois sets the
isomorphism is given by the conjugating group element.
Thus the automorphism of $BG_k$ induced by $\sigma$ is
naturally isomorphic to the identity.

Consider the absolute scheme automorphism $f_\sigma$ of
$\mathbb G_{m,k}$ defined by
\[
 f_\sigma^*:k[T,T^{-1}]\longrightarrow k[T,T^{-1}],
       \qquad c\longmapsto\sigma(c),\quad T\longmapsto T.
\]
Its full small-topos morphism, with the preceding structural
isomorphism composed in, is a relative morphism over $BG_k$.
It is pinned because the underlying scheme map is an isomorphism.
Its map on $k$-points is $a\mapsto\sigma^{-1}(a)$, as is
seen by pulling back the maximal ideal $(T-b)$.
The point-source relative structure at a rational point is
unique up to relative isomorphism by
\citet[Proposition~3.6]{F_CHW2024}, so choosing the inner
conjugator does not change these rational values.
An algebraic $k$-map inducing this relative morphism would be
$cT^m$. Its values at $1$ and $2$ force $c=1$, $m=1$;
a point moved by $\sigma$ contradicts that conclusion.
Hence the asserted homomorphism is injective.

If $\tau\in\operatorname{Aut}(\bar k)$ commutes with $G_E$,
then it preserves their fixed field $E$. Its restriction to
$E$ acts trivially in $\operatorname{Out}(G_E)$, so injectivity
makes that restriction the identity. Thus
$\tau\in G_E\cap C_{\operatorname{Aut}(\bar k)}(G_E)
=Z(G_E)=1$ by F.39.
Conversely, the stated centralizer condition kills $Z(G_E)$
immediately. For an inner-trivial automorphism of $E$, adjust
an extension to $\bar k$ by the inverse inner conjugator.
The adjusted extension centralizes $G_E$, hence is the identity,
so the original field automorphism is the identity.
This proves the equivalence of the two necessary conditions.
\end{proof}

\begin{claim}{F.41}{Proved}
No characteristic-zero field with finite absolute Galois group
satisfies closed-point-preserving \emph{\'etale} reconstruction.
In particular algebraically closed and real-closed fields fail.
Moreover, for every characteristic-zero field $F$, the
Laurent-series field $F((u))$ fails. In this last case there is
an explicit nonalgebraic, nonconstant, pinned relative geometric
endomorphism of $(\mathbb G_{m,F((u))})_{\mathrm{et}}$.
\end{claim}
\begin{proof}
An algebraically closed characteristic-zero field $K$ has a
nonidentity automorphism. Choose $\bar\Q\subset K$, extend
the nontrivial automorphism of $\Q(\sqrt2)$ to $\bar\Q$,
fix a transcendence basis of $K/\bar\Q$, and extend further
to the algebraic closure $K$ of the resulting rational-function
field. This is an automorphism of $K$, not merely an embedding,
because its image is algebraically closed and $K$ is algebraic
over that image. Since $G_K=1$, F.40 excludes $K$.
If $G_k$ is finite, its separable closure is a finite algebraically
closed extension of $k$ in characteristic zero; F.38 would
force this excluded field to reconstruct. For a real-closed
field the extension $k(i)$ is algebraically closed of degree two.
This conclusion includes real-closed fields with trivial field
automorphism group; it also follows directly from F.39 since
their absolute Galois group has order two.

Now put $K=F((u))$ and choose $\bar F$ and compatible roots
of $u$. The exact Puiseux closure is
\[
 \Omega=\bigcup_{\substack{F'/F\ \mathrm{finite}\\n\ge1}}
                   F'((u^{1/n})),\qquad F'\subset\bar F.
\]
It is an algebraic closure of $K$ by
\citet[Proposition, p.~314]{F_vanDenDriesRibenboim1984}.
The finite coefficient-field restriction is essential;
arbitrary series with coefficients in $\bar F$ need not be
algebraic over $K$.

The substitution $\sigma(u)=u+u^2$, fixing $F$, is an
automorphism of $K$; its inverse is the unique series
\[
 v=\frac{-1+\sqrt{1+4u}}2\in u\Q[[u]],\qquad v+v^2=u.
\]
Extend it to $\tau\in\operatorname{Aut}(\Omega)$ by fixing
$\bar F$ and setting
\[
 \tau(u^{1/n})=u^{1/n}(1+u)^{1/n}.
\]
The second factor is the binomial series with rational
coefficients and constant term one. On each
$F'((u^{1/n}))$ this is substitution of a uniformizer with
leading coefficient one. The formulas and their inverse
substitutions are compatible when $F'$ grows and when the
root degree increases: the constant-term-one roots are unique.
They therefore define the asserted automorphism of $\Omega$.

For $g\in G_K$, write $g(u^{1/n})=\zeta_n u^{1/n}$,
$\zeta_n\in\bar F$. Since $g$ fixes $u$ and $\tau$ fixes
$\bar F$,
\[
 g\tau(u^{1/n})
   =\zeta_nu^{1/n}(1+u)^{1/n}
   =\tau g(u^{1/n}).
\]
The two maps also commute on $\bar F$ and on $K$.
Each field $F'((u^{1/n}))$ equals the finite extension
$K(F',u^{1/n})$, so checking these field generators proves
$g\tau=\tau g$ on all of $\Omega$. This avoids any assertion
that arbitrary field automorphisms preserve infinite series
sums. Thus $\sigma$ induces the identity outer action on
$G_K$, although $\sigma(u)\ne u$.

F.40's construction now supplies the full pinned relative
topos endomorphism. Its underlying scheme map as an absolute
map is an isomorphism, so its point map is genuinely nonconstant
and preserves closed points. On $K$-points it is
$a\mapsto\sigma^{-1}(a)$; the three values $1,2,u$ already
exclude every algebraic $K$-endomorphism of $\mathbb G_m$.
These necessary obstructions and negative families do not
prove a converse to F.37 or classify all fields in F.9.
\end{proof}
